% ch4.tex -- Module 7, Chapter 4: Light and special relativity
\chapter{Light and Special Relativity}\label{ch4:chap}
\covers{M7.IQ4.01, M7.IQ4.02, M7.IQ4.03, M7.IQ4.04}

\inquiry{How does the behaviour of light affect concepts of time, space and matter?}

\begin{goals}
  \item state Einstein's two postulates and evaluate the evidence for them
  \item explain simultaneity, time dilation and length contraction
  \item calculate time dilation, length contraction and relativistic momentum
  \item explain why nothing with mass can reach the speed of light
  \item use $E = mc^2$ for the Sun, annihilation and burning fuel
  \item describe the key experiments and judge how well they support relativity
\end{goals}

\begin{syllabus}
  \item M7.IQ4.01: evaluate the evidence for and against Einstein's two postulates
  \item M7.IQ4.02: evidence for time dilation and length contraction, with calculations
  \item M7.IQ4.03: relativistic momentum and the speed limit $c$
  \item M7.IQ4.04: $E = mc^2$ for the Sun, annihilation and burning fuel
\end{syllabus}

%% ==================================================================
\section{Frames of reference and Einstein's postulates}\label{ch4:sec:postulates}
\covers{M7.IQ4.01}

\begin{bigidea}
Motion is always measured from a viewpoint. Einstein said that no steady viewpoint is special,
and that everyone measures the same speed for light in a vacuum. The rest of this chapter
follows from these two ideas.
\end{bigidea}

\begin{keywords}
  \kw{Frame of reference}{a viewpoint, with its own rulers and clocks.}
  \kw{Inertial frame}{a frame that is not accelerating: at rest or at constant velocity.}
  \kw{Postulate}{a starting statement that a theory takes as true.}
  \kw{Aether}{an invisible substance once thought to carry light. It does not exist.}
\end{keywords}

A \term{frame of reference} is a viewpoint, with its own rulers and clocks, from which we
measure positions, times and motion. An \term{observer} measures in one frame. An \term{event}
happens at one place and one instant, such as a flash of light.

\begin{sidebyside}[0.4]
  \centering
  \begin{tikzpicture}[x=0.62cm, y=0.62cm]
    \foreach \sy/\ttl in {0/{(a) constant velocity: \textbf{inertial}},
                          -3.9/{(b) speeding up: \textbf{non-inertial}}} {
      \begin{scope}[yshift=\sy cm]
        \draw[object, fill=hscbrand!5] (0,0) rectangle (5.4,2.2);
        \fill[hscink] (0.9,-0.17) circle (0.17) (4.5,-0.17) circle (0.17);
        \draw[hscink, line width=0.8pt] (-0.3,-0.34) -- (5.7,-0.34);
        \node[diagramlabel, font=\scriptsize] at (2.7,-0.8) {\ttl};
      \end{scope}
    }
    \draw[hscink, line width=0.7pt] (1.0,2.2) -- (1.0,1.05);
    \fill[hscforce] (1.0,1.0) circle (0.11);
    \draw[hscvel, line width=1pt, ->] (3.05,0.75) -- (3.05,1.8);
    \draw[hscvel, line width=1pt, ->] (3.35,1.8) -- (3.35,0.75);
    \fill[object] (3.2,1.95) circle (0.1);
    \fill[hscink] (3.2,0.6) circle (0.07);
    \node[diagramlabel, font=\scriptsize, anchor=north] at (1.4,0.86) {pendulum};
    \node[diagramlabel, font=\scriptsize, anchor=west] at (3.42,1.95) {ball};
    \node[diagramlabel, font=\scriptsize, anchor=west] at (3.4,0.55) {hand};
    \draw[velocity] (1.6,2.6) -- (2.9,2.6) node[right, text=hscvel, font=\scriptsize] {$v$ constant};
    \begin{scope}[yshift=-3.9cm]
      \draw[guide] (1.2,2.2) -- (1.2,0.9);
      \draw[hscink, line width=0.7pt] (1.2,2.2) -- ++(250:1.15) coordinate (bob);
      \fill[hscforce] (bob) circle (0.11);
      \draw[hscvel, line width=1pt, ->] (3.5,0.75) .. controls (3.45,2.45) and (2.75,2.45) .. (2.6,0.75);
      \fill[hscink] (3.5,0.6) circle (0.07);
      \node[diagramlabel, font=\scriptsize, anchor=east] at (2.5,0.55) {lands here};
      \node[diagramlabel, font=\scriptsize, anchor=west] at (3.62,0.55) {hand};
      \draw[accel] (1.6,2.6) -- (2.9,2.6) node[right, text=hscacc, font=\scriptsize] {$a$};
    \end{scope}
  \end{tikzpicture}
  \captionof{figure}{A pendulum and a ball tossed straight up inside a closed
    carriage.}\label{ch4:fig:frames}
\nextside
  An \term{inertial frame of reference} is not accelerating. It is at rest or moves at a
  constant velocity (constant speed in a straight line). Newton's first law works in it.
  A \term{non-inertial frame} is accelerating: it speeds up, slows down or turns. A car going
  round a roundabout at a constant speed is non-inertial.

  \Cref{ch4:fig:frames} shows a test. (a) In an inertial carriage the pendulum hangs straight
  down and the ball lands back in the hand. (b) When the carriage speeds up, the pendulum tilts
  back and the ball lands behind the hand. Most HSC questions treat the Earth's surface as
  inertial.
\end{sidebyside}

\begin{refresher}{Relative velocity and Galilean relativity}
In everyday life, velocities add. You walk forward at \qty{2}{m.s^{-1}} in a train moving at
\qty{20}{m.s^{-1}}, so the platform measures $20 + 2 = \qty{22}{m.s^{-1}}$.

\term{Galilean relativity} says that the laws of motion are the same in every inertial frame.
On a smooth train with the blinds down, a ball you drop falls straight down. It does this whether
the train moves steadily or stands at the station. No experiment with moving objects inside can
tell the difference. (Einstein's first postulate, below, extends this to \emph{all} experiments,
including ones with light.)
\end{refresher}

Light was a puzzle. Maxwell's equations (\cref{ch1:chap}) give its speed,
$c = \qty{3.00e8}{m.s^{-1}}$, but do not say \emph{relative to what}.

Scientists in the 1800s thought light needed a \term{medium}: a substance that a wave travels
through, as sound travels through air. They imagined an invisible medium called the
\term{luminiferous aether}. The Earth moves through space at about \qty{30}{km.s^{-1}}. So, if
the aether existed, light should travel at slightly different speeds in different directions.
In 1887 Michelson and Morley found no difference (\cref{ch4:sec:experiments}).

\subsection{Einstein's two postulates (1905)}
A \term{postulate} is a starting statement that a theory takes as true. Einstein's
\term{special relativity} is the theory of how space and time behave for inertial observers.
It rests on two postulates:
\begin{enumerate}
  \item \textbf{All inertial frames of reference are equivalent.} The laws of physics, including
        those of light, are the same in every inertial frame. No experiment can show which
        inertial frame is ``really'' at rest.
  \item \textbf{The speed of light in a vacuum is an absolute constant.} Every inertial observer
        measures $c = \qty{3.00e8}{m.s^{-1}}$, however fast the source or the observer moves.
\end{enumerate}
The two fit together. Maxwell's equations are laws of physics, so by postulate 1 they give the
same $c$ in every inertial frame. Postulate 2 means that velocities do \emph{not} add for light.
A car's headlight beam travels at $c$, not $c + v$, for every observer.

Speed is distance divided by time. Two observers moving relative to each other measure the same
speed for the same light beam. So they must disagree about the distance, the time, or both
(\cref{ch4:sec:time}).

\Cref{ch4:tab:evidence} sums up the evidence that confirms or denies the postulates.

\begin{table}
  \caption{Evidence about Einstein's postulates.}\label{ch4:tab:evidence}
  \centering\small
  \begin{tabular}{@{}L{36mm}L{78mm}L{20mm}@{}}
    \toprule
    \colhead{Evidence} & \colhead{What was found} & \colhead{Supports}\\
    \midrule
    Maxwell's equations (1860s) & one speed $c$ for light, with no frame mentioned & both\\
    Michelson--Morley (1887) & no change in light's speed in any direction, so the Earth's
      motion could not be detected: a \term{null result} (no effect found) & both\\
    No test detects absolute motion & Michelson--Morley and every later test find the same laws
      of physics in labs on the moving Earth, whatever its direction & postulate 1\\
    Fast pions (short-lived particles), CERN 1964 & light from a source at $0.99975c$ still
      travels at $c$ (\cref{ch4:we:pion}) & postulate 2\\
    Modern laser tests & $c$ is the same in all directions to about 1 part in $10^{17}$ & postulate 2\\
    Muons, atomic clocks, accelerators, GPS (satellite navigation) & every prediction of
      relativity confirmed (\cref{ch4:sec:experiments}) & both\\
    OPERA (2011): a claim \emph{against} relativity's speed limit & neutrinos (tiny, almost
      massless particles) seemed to beat light by \qty{60}{ns}; the cause was a loose
      fibre-optic cable, and repeat tests (ICARUS) found no effect & neither\\
    \bottomrule
  \end{tabular}
\end{table}

Judgement: the evidence strongly supports both postulates. Many independent tests over 100
years agree, and no result against them has survived checking.

One limitation: Michelson and Morley used a light source at rest in their lab. On its own, their
result does not show that a \emph{moving} source gives the same $c$. The 1964 pion test does
(\cref{ch4:we:pion}).

\begin{workedexample}{Speeds measured from the shore}\label{ch4:we:car}
A speedboat moves in a straight line at a constant \qty{15}{m.s^{-1}}. A passenger throws a ball
forward at \qty{12}{m.s^{-1}} relative to the boat. Then she switches on a forward-pointing laser
pointer. Find the speed of (a) the ball and (b) the light, measured by an observer on the shore.
\given{boat \qty{15}{m.s^{-1}}; ball \qty{12}{m.s^{-1}} relative to the boat; light $c$ relative
  to the boat}
\find{the speeds of the ball and of the light relative to the shore}
\formula{everyday velocity addition for the ball; postulate 2 for the light}
\substitute{\eqlines{v_\mathrm{ball} &= \qty{15}{m.s^{-1}} + \qty{12}{m.s^{-1}} = \qty{27}{m.s^{-1}}\\
  v_\mathrm{light} &= c = \qty{3.00e8}{m.s^{-1}} \quad \text{(not $c + \qty{15}{m.s^{-1}}$)}}}
\answer{(a) \final{\qty{27}{m.s^{-1}}} \quad (b) \final{\qty{3.00e8}{m.s^{-1}}}, the same as
  the passenger measures}
\checkstep{addition would give \qty{300000015}{m.s^{-1}}, only 5 parts in 100 million more
  than $c$. A slow source cannot test postulate 2. A very fast one can (\cref{ch4:we:pion}).}
\end{workedexample}

\begin{workedexample}{Light from a source moving at nearly $c$}\label{ch4:we:pion}
Pions are short-lived particles made in accelerators. In an experiment like one at CERN in
1964, pions moved at $0.99975c$. They gave off gamma rays (a type of light) in their direction of
motion. Detectors \qty{40.0}{m} apart timed the gamma rays. Predict the travel time (a) if the
source's speed adds to $c$ (b) using postulate 2. (c) The measured time is \qty{133}{ns}. Which
prediction does it support?
\given{$v = 0.99975c$; distance \qty{40.0}{m}; measured time \qty{133}{ns}
  ($\qty{1}{ns} = \qty{e-9}{s}$)}
\find{the two predicted travel times}
\formula{time $= \dfrac{\text{distance}}{\text{speed}}$}
\substitute{\eqlines{\text{(a) speed} &= (1 + 0.99975) \times \qty{3.00e8}{m.s^{-1}} = \qty{5.99925e8}{m.s^{-1}}\\
  t &= \frac{\qty{40.0}{m}}{\qty{5.99925e8}{m.s^{-1}}} = \qty{6.67e-8}{s} = \qty{66.7}{ns}\\
  \text{(b) } t &= \frac{\qty{40.0}{m}}{\qty{3.00e8}{m.s^{-1}}} = \qty{1.33e-7}{s} = \qty{133}{ns}}}
\answer{(a) \final{\qty{66.7}{ns}} \quad (b) \final{\qty{133}{ns}} \quad (c) the measurement
  matches (b), so it supports postulate 2: light's speed does not depend on its source's speed.}
\checkstep{\cmark\ the predictions differ by a factor of 2. Timers measure to better than
  \qty{1}{ns}, so the test can tell them apart easily.}
\end{workedexample}

\begin{mistakes}
\begin{itemize}
  \item \xmark\ ``Inertial means constant speed.'' \cmark\ Constant \emph{velocity}: moving in a
        circle at constant speed is accelerating.
  \item \xmark\ ``Light's speed is always the same.'' \cmark\ Only \emph{in a vacuum}. Light slows
        in glass and water ($n_\mathrm{x} = \dfrac{c}{v_\mathrm{x}}$).
  \item \xmark\ ``Michelson and Morley proved relativity.'' \cmark\ Their null result
        \emph{supports} both postulates. Evidence supports a theory; it does not prove it.
\end{itemize}
\end{mistakes}

\begin{examtip}
``Evaluate the evidence for the second postulate'' is worth 4--5 marks. Full marks need:
\begin{enumerate}
  \item the postulate stated exactly (``in a vacuum'', ``every inertial observer'');
  \item two or more pieces of evidence and what each found;
  \item a link from each piece of evidence to the postulate;
  \item a clear judgement (``strongly supported: many independent tests agree and none
        disagrees'') plus one specific limitation (for example, Michelson--Morley used a source
        at rest).
\end{enumerate}
For the \emph{first} postulate, use the same plan. Evidence: no experiment, from
Michelson--Morley to modern laser tests, can detect the Earth's own motion. Also, the
predictions built on it (muons, atomic clocks, GPS) are confirmed.
\end{examtip}

\begin{quickcheck}
  \item Which of these is an inertial frame? Give a reason. (A) a cyclist riding along a
        straight, flat road at a steady \qty{6}{m.s^{-1}} (B) a bus turning a corner at a steady
        speed (C) a lift slowing down as it reaches the top floor (D) a Ferris wheel seat moving
        at a steady speed
  \item State Einstein's two postulates of special relativity.
  \item Light travels at \qty{2.25e8}{m.s^{-1}} in water. A student says this disproves the
        second postulate. Explain why the student is wrong.
\end{quickcheck}

%% ==================================================================
\section{Time dilation and length contraction}\label{ch4:sec:time}
\covers{M7.IQ4.02}

\begin{bigidea}
If everyone measures the same speed for light, observers moving relative to each other must
disagree about time and length. A moving clock is measured to run slow. A moving object is
measured to be shorter along its motion.
\end{bigidea}

\begin{keywords}
  \kw{Thought experiment}{an imagined experiment, worked through with logic.}
  \kw{Time dilation}{a moving clock is measured to run slow.}
  \kw{Proper time ($t_0$)}{time between two events, measured by a clock present at both.}
  \kw{Length contraction}{a moving object is measured shorter along its motion.}
  \kw{Proper length ($l_0$)}{length measured in the frame where the object is at rest.}
\end{keywords}

A \term{thought experiment} is an experiment done in your head: you imagine a situation, apply
the postulates and follow the logic. Einstein could not ride a train at nearly the speed of
light, but he could imagine one. A thought experiment is \emph{not} evidence. It shows what the
theory predicts, and real experiments then test it (\cref{ch4:sec:experiments}).

\subsection{Thought experiment 1: lightning and a fast train}
Two events are \term{simultaneous} if they happen at the same time. In \cref{ch4:fig:lightning}
lightning strikes both ends of a very fast train and leaves marks A and B on the track. Olivia
stands on the platform halfway between the marks. Peter sits in the middle of the train.

\begin{figure}
  \centering
  \begin{tikzpicture}[x=0.95cm, y=0.85cm]
    \foreach \yy/\xx/\lab in {0/0/{(1) Lightning strikes both ends of the train.},
        -2.7/1.333/{(2) The front flash reaches Peter first. Both flashes are equally far
          from Olivia.}} {
      \node[diagramlabel, font=\footnotesize, anchor=west] at (0.5,\yy+1.6) {\lab};
      \draw[hscink, line width=0.8pt] (-0.7,\yy) -- (10.2,\yy);
      \foreach \mx/\ml in {0/A, 8/B} {
        \draw[hscmuted, line width=1pt] (\mx-0.12,\yy-0.12) -- (\mx+0.12,\yy+0.12)
          (\mx-0.12,\yy+0.12) -- (\mx+0.12,\yy-0.12);
        \node[diagramlabel, font=\scriptsize, text=hscmuted] at (\mx,\yy-0.36) {\ml};
      }
      \draw[object] (\xx,\yy+0.12) rectangle (\xx+8,\yy+0.9);
      \node[diagramlabel, font=\scriptsize, anchor=west] at (\xx+0.12,\yy+0.51) {rear};
      \node[diagramlabel, font=\scriptsize, anchor=east] at (\xx+7.88,\yy+0.51) {front};
      \fill[hscvel] (\xx+4,\yy+0.42) circle (0.09);
      \node[diagramlabel, font=\scriptsize, text=hscvel, anchor=east] at (\xx+3.85,\yy+0.45) {Peter};
      \fill[hscforce] (4,\yy-0.42) circle (0.09);
      \node[diagramlabel, font=\scriptsize, text=hscforce, anchor=west] at (4.1,\yy-0.48) {Olivia};
    }
    \foreach \bx in {0,8} {
      \draw[hscacc, line width=1.1pt] (\bx+0.15,1.35) -- (\bx-0.05,1.12) -- (\bx+0.12,1.09) -- (\bx-0.02,0.92);
    }
    \draw[velocity] (8.3,0.5) -- (9.3,0.5) node[right, text=hscvel] {$v$};
    \draw[hscray, line width=1.4pt] (2.667,-3.3) -- (2.667,-1.7);
    \draw[hscray, line width=1pt, ->] (2.667,-1.65) -- ++(0.55,0);
    \draw[hscray, line width=1.4pt] (5.333,-3.3) -- (5.333,-1.7);
    \draw[hscray, line width=1pt, ->] (5.333,-1.65) -- ++(-0.55,0);
    \node[diagramlabel, font=\scriptsize, text=hscray, anchor=west, align=left] at (9.4,-2.25)
      {flash\\[-2pt]fronts};
  \end{tikzpicture}
  \caption{Lightning and a fast train, drawn in Olivia's (the platform's) frame. The orange
    lines are the fronts of the two flashes.}\label{ch4:fig:lightning}
\end{figure}

\begin{itemize}
  \item For Olivia the strikes are simultaneous. She is halfway between the marks and both
        flashes travel at $c$, so they reach her together.
  \item The train carries Peter towards the front flash, so it reaches him first.
  \item All observers agree about what happens at one place. So the front flash reaches Peter
        first in \emph{every} frame.
  \item In Peter's own frame he is at rest in the middle of the train. Both flashes travel
        at $c$ (postulate 2) over equal distances. Yet the front flash arrives first. So, for
        Peter, the front strike \emph{really happened first}. By postulate 1 his frame is just
        as valid as Olivia's.
\end{itemize}
This is the \term{relativity of simultaneity}. Two events at different places can be
simultaneous in one inertial frame. In a second frame, moving along the line between them, they
happen at different times. This is not a trick of what each person sees. Each observer has
already allowed for the time the light takes to travel.

\subsection{Thought experiment 2: the light clock}
\begin{sidebyside}[0.5]
  \centering
  \begin{tikzpicture}[x=0.82cm, y=0.82cm]
    \draw[hscink, line width=2pt] (-0.5,0) -- (0.5,0) (-0.5,2) -- (0.5,2);
    \node[diagramlabel, font=\scriptsize, anchor=south] at (0,2.08) {mirror};
    \node[diagramlabel, font=\scriptsize, anchor=north] at (0,-0.08) {mirror};
    \draw[ray, ->] (-0.13,0.06) -- (-0.13,1.94);
    \draw[ray, ->] (0.13,1.94) -- (0.13,0.06);
    \draw[<->, hscmuted, line width=0.6pt] (0.8,0) -- node[right, text=hscink] {$d$} (0.8,2);
    \node[diagramlabel, font=\footnotesize] at (0.1,-1.0) {(a) at rest};
    \begin{scope}[xshift=2.0cm]
      \foreach \x/\op in {0/25, 2.2/50, 4.4/100}
        \draw[hscink!\op, line width=2pt] (\x-0.5,0) -- (\x+0.5,0) (\x-0.5,2) -- (\x+0.5,2);
      \draw[ray] (0,0.04) -- (2.2,1.96);
      \draw[ray] (2.2,1.96) -- (4.4,0.04);
      \draw[guide] (0,-0.08) -- (2.2,-0.08) -- (2.2,1.96);
      \node[diagramlabel, font=\footnotesize, below] at (1.1,-0.12) {$vt/2$};
      \node[diagramlabel, font=\footnotesize, right] at (2.22,1.1) {$d$};
      \node[diagramlabel, font=\footnotesize, above left, text=hscray] at (1.0,1.0) {$ct/2$};
      \draw[velocity] (0.6,2.45) -- (1.8,2.45) node[right, text=hscvel] {$v$};
      \node[diagramlabel, font=\footnotesize] at (2.2,-1.0) {(b) moving at speed $v$};
    \end{scope}
  \end{tikzpicture}
  \captionof{figure}{A light clock in its own frame (a) and in a frame where it moves
    (b).}\label{ch4:fig:lightclock}
\nextside
  A \term{light clock} is a pulse of light bouncing between two mirrors a distance $d$ apart.
  One tick is one trip up and back (\cref{ch4:fig:lightclock}).
  \begin{itemize}
    \item (a) In its own frame, the light travels $2d$, so one tick takes $t_0 = \dfrac{2d}{c}$.
    \item (b) For an observer for whom the clock moves, the mirrors move sideways while the
          light travels. So the light must go on a slant to reach the top mirror. The slanted
          path is longer, but the light still travels at $c$ (postulate 2). So each tick takes
          longer.
  \end{itemize}
\end{sidebyside}

A moving clock ticks more slowly. This is \term{time dilation}. Pythagoras' theorem for the
triangle in \cref{ch4:fig:lightclock}(b) gives the formula. The algebra uses three skills.

\begin{toolkit}{Squares, common factors and square roots}
\begin{itemize}
  \item Squaring a fraction squares the top and the bottom:
        $\left(\dfrac{ct}{2}\right)^2 = \dfrac{c^2t^2}{4}$.
  \item Taking out a common factor: $c^2t^2 - v^2t^2 = t^2\left(c^2 - v^2\right)$. Multiply the
        bracket out again to check.
  \item The square root of a product is the product of the square roots:
        $\sqrt{c^2 - v^2} = \sqrt{c^2\left(1 - \dfrac{v^2}{c^2}\right)}
        = c\sqrt{1 - \dfrac{v^2}{c^2}}$.
\end{itemize}
\end{toolkit}

Each line below does one step.
\begin{align}
  \left(\frac{ct}{2}\right)^2 &= d^2 + \left(\frac{vt}{2}\right)^2
    && \text{Pythagoras: slant$^2$ = height$^2$ + base$^2$}\notag\\
  \frac{c^2t^2}{4} &= d^2 + \frac{v^2t^2}{4} && \text{square each bracket}\notag\\
  c^2t^2 &= 4d^2 + v^2t^2 && \text{multiply every term by 4}\notag\\
  c^2t^2 - v^2t^2 &= 4d^2 && \text{subtract $v^2t^2$ from both sides}\notag\\
  t^2\left(c^2 - v^2\right) &= 4d^2 && \text{take out the common factor $t^2$}\notag\\
  t^2 &= \frac{4d^2}{c^2 - v^2} && \text{divide both sides by $(c^2 - v^2)$}\notag\\
  t &= \frac{2d}{\sqrt{c^2 - v^2}} && \text{square root of both sides}\notag\\
  t &= \frac{2d}{c\sqrt{\left(1 - \dfrac{v^2}{c^2}\right)}}
    && \text{take $c^2$ out of the root (toolkit)}\notag\\
  t &= \frac{t_0}{\sqrt{\left(1 - \dfrac{v^2}{c^2}\right)}}
    && \text{replace $\dfrac{2d}{c}$ by $t_0$}\label{ch4:eq:dilation}
\end{align}

\subsection{Proper time and proper length}
The \term{proper time} $t_0$ is the time between two events measured by a clock that is present
at \emph{both} events. Observers moving relative to that clock measure a longer time $t$. Ask:
\emph{whose clock is at both events?} For a muon's life it is the muon's own clock. For a space
trip it is the ship's clock, which is at the departure and the arrival.

\begin{sidebyside}[0.46]
  \centering
  \begin{tikzpicture}[x=0.8cm, y=0.8cm]
    \draw[object] (0,0) -- (3.2,0) -- (3.7,0.4) -- (3.2,0.8) -- (0,0.8) -- cycle;
    \draw[<->, hscmuted, line width=0.6pt] (0,-0.25) -- node[below, text=hscink, font=\footnotesize] {$l_0$} (3.7,-0.25);
    \node[diagramlabel, font=\footnotesize, anchor=east] at (-0.1,0.4) {at rest};
    \begin{scope}[yshift=-1.75cm]
      \draw[object] (0,0) -- (1.92,0) -- (2.22,0.4) -- (1.92,0.8) -- (0,0.8) -- cycle;
      \draw[<->, hscmuted, line width=0.6pt] (0,-0.25) -- node[below, text=hscink, font=\footnotesize] {$l = 0.60\,l_0$} (2.22,-0.25);
      \draw[velocity] (2.45,0.4) -- (3.55,0.4) node[right, text=hscvel, font=\footnotesize] {$0.80c$};
      \node[diagramlabel, font=\footnotesize, anchor=east] at (-0.1,0.4) {moving};
    \end{scope}
  \end{tikzpicture}
  \captionof{figure}{Only the length along the motion shrinks; the height stays the
    same.}\label{ch4:fig:ship}
\nextside
  \term{Length contraction} means that a moving object is measured to be shorter along its
  direction of motion, but not across it (\cref{ch4:fig:ship}). The \term{proper length} $l_0$
  is measured by an observer for whom the object is at rest. Now think of a distance between two
  places, such as the Earth and a star. Its proper length is measured in the frame where both
  places are at rest. Everyone else measures a shorter length $l$.
\end{sidebyside}

Length contraction follows from time dilation. A ship flies at speed $v$ from the Earth to a
star. The Earth and the star are at rest relative to each other. So the Earth measures the
proper distance $l_0$ between them.
\begin{align}
  t &= \frac{l_0}{v} && \text{Earth frame: time $=$ distance $\div$ speed}\notag\\
  t_0 &= t\sqrt{\left(1 - \frac{v^2}{c^2}\right)}
    && \text{ship's clock is at both events; multiply \cref{ch4:eq:dilation} by the root}\notag\\
  l &= v\,t_0 && \text{ship frame: the star comes towards the ship at $v$ for time $t_0$}\notag\\
  l &= v \times \frac{l_0}{v}\sqrt{\left(1 - \frac{v^2}{c^2}\right)}
    && \text{put in $t_0$ from line 2 and $t$ from line 1}\notag\\
  l &= l_0\sqrt{\left(1 - \frac{v^2}{c^2}\right)}
    && \text{the $v$ on top and the $v$ below cancel}\label{ch4:eq:contraction}
\end{align}
Both effects work both ways. An astronaut measures Earth's clocks running slow. Passengers on a
fast train measure the platform as contracted. Your own clock and your own ruler always seem
normal to you.

\begin{toolkit}[breakable]{Working with speeds near $c$}
\begin{itemize}
  \item $v = 0.80c$ means $\dfrac{v}{c} = 0.80$, so $\dfrac{v^2}{c^2} = 0.80^2 = 0.64$. The $c$
        cancels. If $v$ is in \si{m.s^{-1}}, divide by \num{3.00e8} first.
  \item Order: square, subtract from 1, square root: $\sqrt{1 - 0.64} = \sqrt{0.36} = 0.60$.
        The factor $\dfrac{1}{\sqrt{1 - v^2/c^2}}$ is 1.005 at $0.10c$, 2.29 at $0.90c$ and
        7.09 at $0.99c$. Some books and answers call it the \term{Lorentz factor} $\gamma$. You
        do not need this name: the HSC sheet always writes the full root.
  \item To find $v$ from $t_0$ and $t$, start from $t_0 = t\sqrt{1 - v^2/c^2}$:
        \begin{enumerate}
          \item divide both sides by $t$: $\dfrac{t_0}{t} = \sqrt{1 - \dfrac{v^2}{c^2}}$;
          \item square both sides: $\left(\dfrac{t_0}{t}\right)^2 = 1 - \dfrac{v^2}{c^2}$;
          \item swap the two terms over: $\dfrac{v^2}{c^2} = 1 - \left(\dfrac{t_0}{t}\right)^2$;
          \item square root, then multiply by $c$: $v = c\sqrt{1 - (t_0/t)^2}$.
        \end{enumerate}
        For lengths, use the same steps with $\dfrac{l}{l_0}$.
  \item A \term{light-year} (ly) is the distance light travels in one year (about
        \qty{9.46e15}{m}, not on the data sheet). Something moving at $0.90c$ covers 0.90 ly
        each year. So time in years $=$ distance in ly $\div$ $\dfrac{v}{c}$. Example: 4.5 ly
        at $0.90c$ takes $4.5 \div 0.90 = 5.0$ years. Also $\qty{1}{\micro\second} = \qty{e-6}{s}$.
\end{itemize}
\end{toolkit}

\begin{formulacard}{Time dilation}
  \cardformula{t = \frac{t_0}{\sqrt{\left(1 - \dfrac{v^2}{c^2}\right)}}}
  \begin{symbols}
    \sym{$t$}{time measured in a frame where the clock moves}{\si{s}}
    \sym{$t_0$}{proper time: by a clock present at both events}{\si{s}}
    \sym{$v$}{relative speed of the clock and the observer}{\si{m.s^{-1}}}
    \sym{$c$}{speed of light in a vacuum, \qty{3.00e8}{m.s^{-1}}}{\si{m.s^{-1}}}
  \end{symbols}
  \inwords{a moving clock is measured to run slow, so $t$ is longer than $t_0$.}
  \whentouse{the time between the same two events, in two inertial frames.}
  \whennottouse{for lengths; or when you measure from an accelerating frame. Always measure $v$
    from an inertial frame (for aeroplanes, the Earth's centre). A clock moving in a circle at a
    steady speed is fine: use its speed.}
  \rearranged{$t_0 = t\sqrt{\left(1 - \dfrac{v^2}{c^2}\right)}$ \qquad
    $v = c\sqrt{1 - \left(t_0/t\right)^2}$}
  \tensecond{$t_0 = \qty{1.00}{s}$, $v = 0.600c$:
    $t = \dfrac{1.00}{\sqrt{1 - 0.360}} = \dfrac{1.00}{0.800} = \qty{1.25}{s}$.}
  \trap{square $\dfrac{v}{c}$ first: $\sqrt{1 - 0.9^2}$, not $\sqrt{1 - 0.9}$. Check $t > t_0$.}
\end{formulacard}

\begin{formulacard}{Length contraction}
  \cardformula{l = l_0\sqrt{\left(1 - \dfrac{v^2}{c^2}\right)}}
  \begin{symbols}
    \sym{$l$}{length measured in a frame where the object moves}{\si{m}}
    \sym{$l_0$}{proper length: measured where the object is at rest}{\si{m}}
    \sym{$v$}{relative speed of the object and the observer}{\si{m.s^{-1}}}
    \sym{$c$}{speed of light in a vacuum, \qty{3.00e8}{m.s^{-1}}}{\si{m.s^{-1}}}
  \end{symbols}
  \inwords{a moving object is measured shorter along its motion, so $l$ is less than $l_0$.}
  \whentouse{lengths and distances \emph{along} the motion (a ship; Earth to a star; the
    thickness of the atmosphere for a muon).}
  \whennottouse{for heights or widths at right angles to the motion. They do not change.}
  \rearranged{$l_0 = \dfrac{l}{\sqrt{\left(1 - \dfrac{v^2}{c^2}\right)}}$ \qquad
    $v = c\sqrt{1 - \left(l/l_0\right)^2}$}
  \tensecond{$l_0 = \qty{10.0}{m}$, $v = 0.600c$: $l = 10.0 \times 0.800 = \qty{8.00}{m}$.}
  \trap{$l_0$ belongs to the observer for whom the object is \emph{at rest}. If your $l$ is
    longer than $l_0$, you divided instead of multiplied.}
\end{formulacard}

\begin{workedexample}{A clock on a space probe}\label{ch4:we:probe}
A space probe passes the Earth at $0.800c$. A camera on the probe takes photos for
\qty{12.0}{\minute} by the probe's own clock. How long does this last, as measured by mission
control on Earth?
\given{$t_0 = \qty{12.0}{\minute}$ (the probe's clock is at the start and the end); $v = 0.800c$}
\find{the time $t$ measured on Earth}
\formula{$t = \dfrac{t_0}{\sqrt{\left(1 - \dfrac{v^2}{c^2}\right)}}$}
\substitute{\eqlines{\frac{v^2}{c^2} &= 0.800^2 = 0.640 \quad\text{(square first)}\\
  t &= \frac{\qty{12.0}{\minute}}{\sqrt{1 - 0.640}} = \frac{\qty{12.0}{\minute}}{\sqrt{0.360}}
  = \frac{\qty{12.0}{\minute}}{0.600} = \qty{20.0}{\minute}}}
\answer{\final{$t = \qty{20.0}{\minute}$} (3 significant figures, as in the data)}
\checkstep{$t > t_0$ \cmark. Minutes can stay as minutes, because the formula compares two times.}
\end{workedexample}

\begin{workedexample}{How fast is the rocket?}\label{ch4:we:rocket}
A rocket is \qty{60.0}{m} long at rest. A space station measures its length as \qty{45.0}{m}
as it flies past. Find the rocket's speed relative to the station.
\given{$l_0 = \qty{60.0}{m}$ (the rocket is at rest in its own frame); $l = \qty{45.0}{m}$
  (for the station the rocket moves)}
\find{the speed $v$}
\formula{$l = l_0\sqrt{\left(1 - \dfrac{v^2}{c^2}\right)}$, rearranged for $v$ (Maths toolkit)}
\substitute{\eqlines{\sqrt{1 - \frac{v^2}{c^2}} &= \frac{l}{l_0} = \frac{\qty{45.0}{m}}{\qty{60.0}{m}} = 0.750
    \quad\text{(divide both sides by $l_0$)}\\
  0.750^2 &= 0.5625 = 1 - \frac{v^2}{c^2} \quad\text{(square both sides)}\\
  \frac{v^2}{c^2} &= 1 - 0.5625 = 0.4375 \quad\text{(swap the terms over)}\\
  \frac{v}{c} &= \sqrt{0.4375} = 0.661 \quad\text{(square root)}\\
  v &= 0.661 \times \qty{3.00e8}{m.s^{-1}} = \qty{1.98e8}{m.s^{-1}}}}
\answer{\final{$v = 0.661c = \qty{1.98e8}{m.s^{-1}}$} (3 significant figures)}
\checkstep{$v < c$ \cmark. Put it back in: $\qty{60.0}{m} \times \sqrt{1 - 0.661^2} = \qty{45.0}{m}$
  \cmark.}
\end{workedexample}

\subsection{Muons: both effects in action}
\begin{sidebyside}[0.5]
  \centering
  \begin{tikzpicture}[x=0.72cm, y=0.72cm]
    \fill[ground] (0,-0.25) rectangle (2.6,0);
    \draw[hscink, line width=0.8pt] (0,0) -- (2.6,0);
    \draw[hscmuted, dashed, line width=0.6pt] (0,3.6) -- (2.6,3.6);
    \draw[<->, hscmuted, line width=0.6pt] (-0.2,0) -- node[left, text=hscink, font=\footnotesize]
      {$l_0$} (-0.2,3.6);
    \node[diagramlabel, font=\scriptsize, anchor=west] at (2.65,-0.12) {ground};
    \node[diagramlabel, font=\scriptsize, anchor=south, align=center] at (1.3,3.62)
      {top of\\[-2pt]atmosphere};
    \fill[object] (1.0,3.3) circle (0.11);
    \node[diagramlabel, font=\scriptsize, anchor=west] at (1.15,3.3) {muon};
    \draw[velocity] (1.0,3.05) -- (1.0,1.9) node[left, text=hscvel, font=\scriptsize] {$v$};
    \node[diagramlabel, font=\scriptsize, align=left, anchor=west] at (1.15,2.4)
      {its clock\\[-2pt]runs slow};
    \node[diagramlabel, font=\footnotesize] at (1.3,-0.7) {(a) Earth frame};
    \begin{scope}[xshift=2.9cm]
      \fill[ground] (0,1.95) rectangle (2.6,2.2);
      \draw[hscink, line width=0.8pt] (0,2.2) -- (2.6,2.2);
      \node[diagramlabel, font=\scriptsize, anchor=west] at (2.65,2.08) {ground};
      \draw[hscmuted, dashed, line width=0.6pt] (0,3.2) -- (2.6,3.2);
      \node[diagramlabel, font=\scriptsize, anchor=south, align=center] at (1.3,3.22)
        {top of\\[-2pt]atmosphere};
      \draw[<->, hscmuted, line width=0.6pt] (-0.2,2.2) -- node[left, text=hscink, font=\footnotesize]
        {$l$} (-0.2,3.2);
      \fill[object] (1.0,2.95) circle (0.11);
      \node[diagramlabel, font=\scriptsize, anchor=west] at (1.15,2.92) {muon at rest};
      \draw[velocity] (1.8,0.55) -- (1.8,1.75);
      \node[diagramlabel, font=\scriptsize, align=right, anchor=east] at (1.65,1.1) {Earth\\[-2pt]rushes up};
      \node[diagramlabel, font=\footnotesize] at (1.3,-0.7) {(b) muon frame};
    \end{scope}
  \end{tikzpicture}
  \captionof{figure}{One muon trip in two frames (not to scale). (a) The atmosphere has its
    proper thickness $l_0$. (b) In the muon frame the atmosphere is contracted to
    $l$.}\label{ch4:fig:muons}
\nextside
  Cosmic rays are fast protons from space. They hit air nuclei 10--15 km above the ground. This
  makes short-lived particles called pions, which quickly turn into \term{muons}. A muon is
  like a heavy electron. It is \term{unstable}: after a short time it \term{decays} (breaks up
  into other particles). At rest a muon has a \term{mean lifetime} (average life) of
  \qty{2.2}{\micro\second}. Even at nearly $c$ it would go only about \qty{660}{m} in this time.
  Yet many muons reach the ground (\cref{ch4:fig:muons}).
  \begin{itemize}
    \item (a) Earth frame: the muon's clock runs slow, so it lives long enough.
    \item (b) Muon frame: the atmosphere is contracted, so the trip is short enough.
  \end{itemize}
\end{sidebyside}

\begin{workedexample}{Muons in two frames}\label{ch4:we:muon}
Muons are made \qty{12.0}{km} above the ground (Earth measurement) and travel straight down at
$0.9990c$. Their mean lifetime at rest is \qty{2.20}{\micro\second}. (a) How far would one go in
one lifetime without relativity? (b) Earth frame: compare the trip time with the dilated
lifetime. (c) Muon frame: find the thickness of the atmosphere and the trip time. (d) What do
(b) and (c) mean for the number of muons that reach the ground?
\given{$l_0 = \qty{1.20e4}{m}$ (atmosphere at rest in the Earth frame); $v = 0.9990c =
  \qty{2.997e8}{m.s^{-1}}$; $t_0 = \qty{2.20e-6}{s}$ (the muon's own clock)}
\find{(a) the distance $vt_0$ travelled in one lifetime without relativity \quad (b) the trip
  time and the dilated lifetime $t$ \quad (c) the contracted thickness $l$ and the trip time
  \quad (d) trip time $\div$ lifetime in each frame}
\formula{distance $= vt$; \ $t = \dfrac{t_0}{\sqrt{\left(1 - \dfrac{v^2}{c^2}\right)}}$; \
  $l = l_0\sqrt{\left(1 - \dfrac{v^2}{c^2}\right)}$}
\substitute{\eqlines{\sqrt{1 - 0.9990^2} &= \sqrt{0.001999} = 0.04471 \quad\text{(work this out once)}\\
  \text{(a) } vt_0 &= \qty{2.997e8}{m.s^{-1}} \times \qty{2.20e-6}{s} = \qty{659}{m}\\
  \text{(b) trip} &= \frac{\qty{1.20e4}{m}}{\qty{2.997e8}{m.s^{-1}}} = \qty{40.0}{\micro\second}\\
  t &= \frac{\qty{2.20}{\micro\second}}{0.04471} = \qty{49.2}{\micro\second}\\
  \text{(c) } l &= \qty{1.20e4}{m} \times 0.04471 = \qty{537}{m}\\
  \text{trip} &= \frac{\qty{537}{m}}{\qty{2.997e8}{m.s^{-1}}} = \qty{1.79}{\micro\second}\\
  \text{(d) } \frac{40.0}{49.2} &= 0.81 \qquad \frac{1.79}{2.20} = 0.81}}
\answer{(a) \final{\qty{659}{m}}, far less than \qty{12.0}{km} \quad (b) trip
  \qty{40.0}{\micro\second} $<$ lifetime \qty{49.2}{\micro\second} \quad (c) \final{$l = \qty{537}{m}$};
  trip \qty{1.79}{\micro\second} $<$ \qty{2.20}{\micro\second} \quad (d) both frames agree: the
  trip lasts 0.81 of a mean lifetime. Decay is random, so some muons decay on the way. But a
  large fraction of them (just under half) reach the ground.}
\checkstep{\cmark\ the ratio is the same in both frames. Without relativity, the trip would last
  $40.0 \div 2.20 = 18$ lifetimes, and almost no muons would arrive. Many do arrive.}
\end{workedexample}

\begin{mistakes}
\begin{itemize}
  \item \xmark\ ``Earth is at rest, so Earth time is proper time.'' \cmark\ Proper time belongs to
        the clock present at both events, whatever frame it is in.
  \item \xmark\ Shrinking the height or width. \cmark\ Only the length along the motion changes.
  \item \xmark\ Using dilated time \emph{and} contracted length in one frame. \cmark\ Earth frame:
        dilated time with proper distance. Muon frame: proper time with contracted distance.
  \item \xmark\ ``Passengers measure the platform as longer.'' \cmark\ Each frame measures the
        other as contracted.
  \item \xmark\ ``Every muon lives exactly \qty{2.2}{\micro\second}.'' \cmark\ That is an average.
        Decay is random: some muons live longer, some shorter.
\end{itemize}
\end{mistakes}

\begin{examtip}
In calculations, one mark is often for choosing $t_0$ or $l_0$. Write a short line, such as
``the ship's clock is at both events, so ship time $= t_0$''. Then substitute and check that
$t > t_0$ or $l < l_0$.

``Explain how muons support special relativity'' is worth 3--5 marks. Full marks need:
\begin{enumerate}
  \item the problem: about \qty{660}{m} in one lifetime, but muons are made 10--15 km up;
  \item Earth frame: the muon's lifetime is dilated;
  \item muon frame: the atmosphere is contracted;
  \item both frames agree, and far more muons reach the ground than classical physics predicts.
\end{enumerate}
\end{examtip}

\begin{quickcheck}
  \item A light on a spaceship flashes every \qty{30.0}{s} according to the crew. The ship moves at
        $0.600c$ relative to Earth. Find the time between flashes measured from Earth.
  \item In \cref{ch4:fig:lightning}, suppose the train moves to the \emph{left}. In Peter's frame,
        which strike happens first? Explain briefly.
  \item A shuttle is \qty{120}{m} long at rest and moves at $0.60c$ past a station. What length is
        measured (a) by the station (b) by the shuttle's pilot?
  \item A star is \qty{4.2}{\lightyear} from Earth (Earth frame). A probe flies there at $0.90c$.
        Find (a) the distance measured by the probe (b) the trip time on the probe's clock.
\end{quickcheck}

%% ==================================================================
\section{Relativistic momentum and the speed limit}\label{ch4:sec:momentum}
\covers{M7.IQ4.03}

\begin{bigidea}
At high speed, momentum grows much faster than $m_0v$. As the speed gets close to $c$, the
momentum heads towards infinity. So nothing with mass can ever reach the speed of light.
\end{bigidea}

\begin{keywords}
  \kw{Rest mass ($m_0$)}{the mass measured by an observer for whom the object is at rest.}
  \kw{Relativistic momentum ($p_\mathrm{v}$)}{the correct momentum at any speed, even near $c$.}
  \kw{Particle accelerator}{a machine that speeds up charged particles with electric fields
    and steers them with magnetic fields.}
\end{keywords}

\begin{refresher}{Momentum and impulse}
Momentum is $p = mv$ (unit \si{kg.m.s^{-1}}). A net force changes momentum:
$\Delta\vec{p} = \vec{F}_\mathrm{net}\Delta t$ (on the sheet). A bigger force for a longer time
gives a bigger change. With no outside force, total momentum is conserved.
\end{refresher}

\begin{sidebyside}[0.47]
  \centering
  \begin{tikzpicture}
    \begin{axis}[hsc, width=\linewidth, height=0.78\linewidth,
        xmin=0, xmax=1.08, ymin=0, ymax=7, xtick={0,0.2,...,1.0},
        xticklabel style={/pgf/number format/.cd, fixed, precision=1},
        xlabel={speed $v/c$}, ylabel={momentum $p/(m_0c)$}]
      \addplot[hscseriesA, no marks, line width=1.2pt, domain=0:0.99, samples=200] {x/sqrt(1-x^2)};
      \addplot[hscseriesB, no marks, dashed, line width=1.2pt, domain=0:1.0, samples=2] {x};
      \draw[guide] (axis cs:1,0) -- (axis cs:1,7);
      \node[font=\footnotesize, anchor=east, text=hscseriesA] at (axis cs:0.92,5.2) {$p_\mathrm{v}$};
      \node[font=\footnotesize, anchor=north, inner sep=1pt, text=hscseriesB] at (axis cs:0.8,0.68) {$m_0v$};
      \node[font=\footnotesize, anchor=north, rotate=90] at (axis cs:1.02,3.5) {$v = c$};
    \end{axis}
  \end{tikzpicture}
  \captionof{figure}{Momentum against speed. Momentum is shown in units of $m_0c$, so the
    numbers work for any particle. The dashed line is $p = m_0v$.}\label{ch4:fig:momentum}
\nextside
  Near $c$, $p = mv$ fails. With $p = mv$, momentum would be conserved in a fast collision in
  one inertial frame but not in another. That would break postulate 1. The \term{relativistic momentum} is correct at
  any speed (formula card below). It uses the \term{rest mass} $m_0$, the mass measured when
  the object is at rest.
  \begin{itemize}
    \item At low speed the root is about 1, so $p_\mathrm{v} \approx m_0v$
          (\cref{ch4:fig:momentum}).
    \item Near $c$, $p_\mathrm{v}$ shoots upwards and never reaches the line $v = c$.
  \end{itemize}
\end{sidebyside}

\begin{toolkit}{Dividing by a very small number}
The arrow $\to$ means ``gets closer and closer to''. So $v \to c$ means ``$v$ gets closer and
closer to $c$''. The symbol $\infty$ means infinity: bigger than any number. Dividing by a small
number gives a big answer: $\dfrac{1}{0.1} = 10$, $\dfrac{1}{0.01} = 100$,
$\dfrac{1}{0.001} = 1000$. The root $\sqrt{1 - v^2/c^2}$ is 0.436 at $0.90c$, 0.141 at $0.99c$
and 0.0447 at $0.999c$. As the bottom of the fraction gets closer to 0, $p_\mathrm{v}$ grows
without limit.
\end{toolkit}

\textbf{Why nothing with mass can reach $c$.} As $v \to c$, the root
$\sqrt{1 - v^2/c^2} \to 0$, so $p_\mathrm{v} \to \infty$. But a real force acting for a real time
gives only a \emph{finite} change in momentum ($\Delta p = F_\mathrm{net}\Delta t$). No push can
supply infinite momentum or energy. So an object with mass gets closer and closer to $c$ but
never reaches it. Each extra bit of speed costs more than the last (\cref{ch4:we:proton}).
A photon has no rest mass, so it always travels at $c$.

\textbf{Consequences and applications.} The Large Hadron Collider gives protons about 7000 times
their rest energy (the energy their mass is worth, \cref{ch4:sec:emc2}). Yet they travel just
below $c$, at about $0.999\,999\,99c$.

A magnet bends a moving charged particle into a circle. From Module 6, the magnetic force $qvB$
provides the centripetal force $\dfrac{mv^2}{r}$, so $qvB = \dfrac{mv^2}{r}$. Multiply both
sides by $r$: $qvBr = mv^2$. Divide both sides by $qvB$:
$r = \dfrac{mv}{qB} = \dfrac{p}{qB}$ \hscbadgenotonsheet. Near $c$, use $p = p_\mathrm{v}$. In a \term{synchrotron}, a ring-shaped accelerator, $r$ is fixed. So the
field $B$ must rise in step with $p_\mathrm{v}$, not $m_0v$.

Some older books say that mass increases at high speed (``relativistic mass''). The HSC formula
keeps $m_0$ fixed: it is the momentum that grows.

\begin{formulacard}{Relativistic momentum}
  \cardformula{p_\mathrm{v} = \frac{m_0 v}{\sqrt{\left(1 - \dfrac{v^2}{c^2}\right)}}}
  \begin{symbols}
    \sym{$p_\mathrm{v}$}{momentum of an object moving at speed $v$}{\si{kg.m.s^{-1}}}
    \sym{$m_0$}{rest mass}{\si{kg}}
    \sym{$v$}{speed of the object relative to the observer}{\si{m.s^{-1}}}
    \sym{$c$}{speed of light in a vacuum, \qty{3.00e8}{m.s^{-1}}}{\si{m.s^{-1}}}
  \end{symbols}
  \inwords{momentum is rest mass times velocity, divided by $\sqrt{1 - v^2/c^2}$. Near $c$
    this root is tiny, so the momentum grows without limit.}
  \whentouse{any moving object. Above about $0.1c$, $p = mv$ is no longer accurate, so you
    \emph{must} use this formula.}
  \whennottouse{for a photon (it has no rest mass and always travels at $c$, so the formula
    does not apply); or with a speed measured by a different observer from the one who
    measures $p_\mathrm{v}$.}
  \rearranged{$m_0 = \dfrac{p_\mathrm{v}}{v}\sqrt{\left(1 - \dfrac{v^2}{c^2}\right)}$ \qquad
    $\dfrac{p_\mathrm{v}}{m_0v} = \dfrac{1}{\sqrt{1 - v^2/c^2}}$}
  \tensecond{an electron at $0.600c$:
    $p_\mathrm{v} = \dfrac{9.109\times10^{-31} \times 1.80\times10^{8}}{0.800}
    = \qty{2.05e-22}{kg.m.s^{-1}}$.}
  \trap{on the top line use $v$ in \si{m.s^{-1}} ($0.600c = \qty{1.80e8}{m.s^{-1}}$), not 0.600.
    Rest mass in kg: change tonnes to kg.}
\end{formulacard}

\begin{workedexample}{An electron at 0.90\texorpdfstring{$c$}{c}}\label{ch4:we:electron}
An electron moves at $0.90c$. Calculate its relativistic momentum and compare it with $m_0v$.
\given{$m_0 = \qty{9.109e-31}{kg}$ (data sheet); $v = 0.90 \times \qty{3.00e8}{m.s^{-1}}
  = \qty{2.70e8}{m.s^{-1}}$}
\find{$p_\mathrm{v}$ and $m_0v$}
\formula{$p_\mathrm{v} = \dfrac{m_0 v}{\sqrt{\left(1 - \dfrac{v^2}{c^2}\right)}}$}
\substitute{\eqlines{\sqrt{1 - 0.90^2} &= \sqrt{1 - 0.81} = \sqrt{0.19} = 0.436
    \quad\text{(root first)}\\
  p_\mathrm{v} &= \frac{\qty{9.109e-31}{kg} \times \qty{2.70e8}{m.s^{-1}}}{0.436}\\
    &= \qty{5.64e-22}{kg.m.s^{-1}}\\
  m_0v &= \qty{9.109e-31}{kg} \times \qty{2.70e8}{m.s^{-1}}\\
    &= \qty{2.46e-22}{kg.m.s^{-1}}}}
\answer{\final{$p_\mathrm{v} = \qty{5.6e-22}{kg.m.s^{-1}}$}, which is
  $\dfrac{5.64}{2.46} = 2.3$ times $m_0v$}
\checkstep{\cmark\ the ratio is $\dfrac{1}{0.436} = 2.29$, the same factor as time dilation at $0.90c$.}
\end{workedexample}

\begin{workedexample}{Faster and faster, but never \texorpdfstring{$c$}{c}}\label{ch4:we:proton}
An accelerator raises a proton's speed from $0.90c$ to $0.99c$. (a) Find its momentum at each
speed. (b) Compare the increase in speed with the increase in momentum. (c) Explain why the
proton can never reach $c$.
\given{$m_0 = \qty{1.673e-27}{kg}$; $v_1 = \qty{2.70e8}{m.s^{-1}}$; $v_2 = \qty{2.97e8}{m.s^{-1}}$}
\find{$p_1$, $p_2$ and the ratios}
\formula{$p_\mathrm{v} = \dfrac{m_0 v}{\sqrt{\left(1 - \dfrac{v^2}{c^2}\right)}}$}
\substitute{\eqlines{p_1 &= \frac{\qty{1.673e-27}{kg} \times \qty{2.70e8}{m.s^{-1}}}{\sqrt{1 - 0.90^2}}
    = \frac{\qty{4.517e-19}{kg.m.s^{-1}}}{0.4359}\\ &= \qty{1.04e-18}{kg.m.s^{-1}}\\
  p_2 &= \frac{\qty{1.673e-27}{kg} \times \qty{2.97e8}{m.s^{-1}}}{\sqrt{1 - 0.99^2}}
    = \frac{\qty{4.969e-19}{kg.m.s^{-1}}}{0.1411}\\ &= \qty{3.52e-18}{kg.m.s^{-1}}\\
  \frac{v_2}{v_1} &= \frac{0.99}{0.90} = 1.10\\
  \frac{p_2}{p_1} &= \frac{3.52}{1.04} = 3.4}}
\answer{(a) \final{\qty{1.0e-18}{kg.m.s^{-1}}} and \final{\qty{3.5e-18}{kg.m.s^{-1}}}
  (b) the speed rises by only 10\,\%, but the momentum more than triples (c) at $v = c$ the
  momentum would be infinite, and a finite force for a finite time ($\Delta p =
  F_\mathrm{net}\Delta t$) can never supply that, so $v$ stays below $c$.}
\checkstep{$p_1$ is 2.3 times $m_0v_1$, the same factor as in \cref{ch4:we:electron} at $0.90c$
  \cmark. Both answers are larger than $m_0v$ \cmark. (Going on to $0.999c$, only 0.9\,\%
  faster, multiplies the momentum by another 3.2.)}
\end{workedexample}

%% Added by the Module 7 coverage audit (M7.IQ4.03, applications of relativistic momentum).
\begin{workedexample}{Steering a fast proton}\label{ch4:we:ring}
In a ring accelerator, a proton moves at $0.960c$ at \ang{90} to a magnetic field of
\qty{1.50}{T}. (a) Find its relativistic momentum. (b) Find the radius of its circular path.
(c) Find the radius that $p = m_0v$ would predict. (d) Explain why the ring's magnets must be
designed with $p_\mathrm{v}$.
\given{$m_0 = \qty{1.673e-27}{kg}$; $v = 0.960 \times \qty{3.00e8}{m.s^{-1}} =
  \qty{2.88e8}{m.s^{-1}}$; $q = \qty{1.602e-19}{C}$; $B = \qty{1.50}{T}$}
\find{(a) $p_\mathrm{v}$ \quad (b) $r$ \quad (c) $r$ from $m_0v$ \quad (d) why $p_\mathrm{v}$
  matters}
\formula{$p_\mathrm{v} = \dfrac{m_0 v}{\sqrt{\left(1 - \dfrac{v^2}{c^2}\right)}}$, \ then
  $r = \dfrac{p}{qB}$ (derived above; not on the sheet)}
\substitute{root first: $\sqrt{1 - 0.960^2} = \sqrt{1 - 0.9216} = \sqrt{0.0784} = 0.280$\\
  (a) top line, then divide by the root:\\
  \eqlines{m_0v &= \qty{1.673e-27}{kg} \times \qty{2.88e8}{m.s^{-1}}\\
    &= \qty{4.818e-19}{kg.m.s^{-1}}\\
  p_\mathrm{v} &= \frac{\qty{4.818e-19}{kg.m.s^{-1}}}{0.280} = \qty{1.721e-18}{kg.m.s^{-1}}}\\[5pt]
  (b) bottom line $qB$ first, then $r$:\\
  \eqlines{qB &= \qty{1.602e-19}{C} \times \qty{1.50}{T} = \qty{2.403e-19}{C.T}\\
  r &= \frac{\qty{1.721e-18}{kg.m.s^{-1}}}{\qty{2.403e-19}{C.T}} = \qty{7.161}{m}}\\[5pt]
  (c) the same, with $m_0v$ on top:\\
  \eqlines{r &= \frac{\qty{4.818e-19}{kg.m.s^{-1}}}{\qty{2.403e-19}{C.T}} = \qty{2.005}{m}}}
\answer{(a) \final{$p_\mathrm{v} = \qty{1.72e-18}{kg.m.s^{-1}}$} \quad
  (b) \final{$r = \qty{7.16}{m}$} \quad (c) \final{$r = \qty{2.01}{m}$}, far too small.
  (d) The real path is $\dfrac{7.161}{2.005} = 3.57$ times wider. Magnets set up with $m_0v$
  would be too weak. The protons would curve too little and hit the outer wall of the ring.}
\checkstep{$\dfrac{1}{0.280} = 3.57$: the radius grows by the same factor as the momentum
  \cmark. Unit: $\unit{kg.m.s^{-1}} \div (\unit{C} \times \unit{T}) = \unit{m}$, because
  \qty{1}{T} $=$ \qty{1}{kg.C^{-1}.s^{-1}} \cmark.}
\end{workedexample}

\begin{mistakes}
\begin{itemize}
  \item \xmark\ Using $p = mv$ at $0.9c$. \cmark\ Use $p_\mathrm{v}$; above about $0.1c$ the
        difference matters.
  \item \xmark\ ``Its mass increases, so it can't go faster.'' \cmark\ Say that the momentum and
        energy needed grow without limit as $v$ approaches $c$.
  \item \xmark\ Forgetting that momentum depends on the frame. \cmark\ For someone moving with the
        object, its momentum is zero.
\end{itemize}
\end{mistakes}

\begin{examtip}
``Explain why a particle with mass cannot reach the speed of light'' is worth 3 marks. Full
marks need:
\begin{enumerate}
  \item as $v \to c$, the root $\to 0$, so $p_\mathrm{v} \to \infty$;
  \item a finite force for a finite time (finite energy) gives only a finite momentum, so
        $v < c$;
  \item evidence: accelerator particles get closer and closer to $c$ but never reach it.
\end{enumerate}
\end{examtip}

\begin{quickcheck}
  \item A cargo probe of rest mass \qty{1.5e5}{kg} moves at $0.80c$ past a space station. Find
        its momentum measured (a) by the station (b) by a camera fixed to the probe.
  \item At what speed is an object's relativistic momentum twice $m_0v$?
  \item Explain why the magnets in a ring accelerator must get stronger as the particles speed up,
        even though their speed hardly changes near $c$.
  \item A proton synchrotron has a ring of radius \qty{25.0}{m}. Find the magnetic field needed
        to keep protons on the ring at (a) $0.990c$ (b) $0.999c$. Compare the increase in speed
        with the increase in field.\label{ch4:qc:ring}
\end{quickcheck}

%% ==================================================================
\section{Mass--energy equivalence: \texorpdfstring{$E = mc^2$}{E = mc2}}\label{ch4:sec:emc2}
\covers{M7.IQ4.04}

\begin{bigidea}
Mass is a very concentrated form of energy. When a process releases energy, the products have
slightly less mass than the starting materials. The energy released is the lost mass times $c^2$.
\end{bigidea}

\begin{keywords}
  \kw{Rest energy}{the energy an object has because of its mass, $E = mc^2$.}
  \kw{Mass defect ($\Delta m$)}{mass before minus mass after: the mass that disappears.}
  \kw{Antiparticle}{a particle with the same mass as another but the opposite charge (and
    other opposite properties).}
  \kw{Annihilation}{a particle and its antiparticle meet and become photons.}
\end{keywords}

Special relativity shows that mass and energy are two forms of the same thing. This is
\term{mass--energy equivalence}: any mass $m$ has a \term{rest energy} $E = mc^2$. Since
$c^2 = \qty{9.00e16}{m^2.s^{-2}}$, a tiny mass holds a huge energy. One gram is equivalent to
\qty{9.0e13}{J}. That is about one day's output of a \qty{1000}{MW} power station
($\qty{1.0e9}{W} \times \qty{86400}{s} = \qty{8.6e13}{J}$).

When a process releases energy (heat, light, motion), the products have slightly less rest mass
than the starting materials. The mass lost is the \term{mass defect}
$\Delta m = m_\mathrm{before} - m_\mathrm{after}$, and $\Delta E = \Delta m\,c^2$. A sealed
container only gets lighter when that energy escapes from it. (In Module 8 you will also meet the
mass defect of a \emph{nucleus}. That is the mass of its separate protons and neutrons minus the
mass of the nucleus.) The syllabus names three examples.

\textbf{1. The Sun.} The Sun shines by \term{nuclear fusion}: light \term{nuclei} (the tiny
centres of atoms) joining to make a heavier one. In its core, four hydrogen nuclei join in steps
to make one helium-4 nucleus (2 protons and 2 neutrons). This is the \term{proton--proton
chain} (Module 8). About 0.7\,\% of the mass becomes energy (\cref{ch4:we:sun}).

\textbf{2. Annihilation.} Every particle has an \term{antiparticle}, with the same mass but
the opposite charge. The electron's antiparticle is called the \term{positron} ($\mathrm{e}^+$).
It has the electron's mass but a positive charge. When they meet, both vanish and all their mass
becomes gamma-ray photons: \term{annihilation}. For a pair at rest, \emph{two} photons fly off in
opposite directions (\cref{ch4:fig:annihilation}), because the total momentum must stay zero.
Hospital PET scanners (positron emission tomography, a type of medical scan) detect these photon
pairs.

\begin{figure}
  \centering
  \begin{tikzpicture}[x=0.9cm, y=0.9cm]
    \fill[hscvel!80] (0,0) circle (0.2) node[white, font=\footnotesize\bfseries] {$-$};
    \fill[hscforce!85] (2.4,0) circle (0.2) node[white, font=\footnotesize\bfseries] {$+$};
    \draw[velocity] (0.3,0) -- (0.95,0);
    \draw[velocity] (2.1,0) -- (1.45,0);
    \node[diagramlabel, font=\footnotesize, below=2.5mm] at (0,0) {electron};
    \node[diagramlabel, font=\footnotesize, below=2.5mm] at (2.4,0) {positron};
    \node[diagramlabel, font=\footnotesize, align=center] at (1.2,-1.45) {before:\\ total momentum $= 0$};
    \draw[->, hscmuted, line width=1pt] (3.3,0) -- (4.2,0);
    \begin{scope}[xshift=6.6cm]
      \fill[hscacc] (0,0) circle (0.07);
      \foreach \dir in {1,-1}
        \draw[hscray, line width=1.1pt, decorate, decoration={snake, amplitude=1.1mm,
          segment length=3mm, post length=2.5mm}, ->] (\dir*0.15,0) -- (\dir*2.1,0);
      \node[diagramlabel, font=\footnotesize, above=2mm, text=hscray] at (1.3,0) {photon};
      \node[diagramlabel, font=\footnotesize, above=2mm, text=hscray] at (-1.3,0) {photon};
      \node[diagramlabel, font=\footnotesize, align=center] at (0,-1.45) {after: two equal photons
        in opposite directions,\\ so total momentum $= 0$};
    \end{scope}
  \end{tikzpicture}
  \caption{Electron--positron annihilation gives two photons, not one.}\label{ch4:fig:annihilation}
\end{figure}

\textbf{3. Burning fuel.} Chemical reactions also lose mass, but very little. Burning
\qty{1.00}{kg} of petrol releases about \qty{4.6e7}{J}, so
$\Delta m = \dfrac{\Delta E}{c^2} = \dfrac{\qty{4.6e7}{J}}{\qty{9.00e16}{m^2.s^{-2}}}
= \qty{5.1e-10}{kg}$: about 5 parts in 10 billion of the fuel's mass. This is far too small to
measure. We know it only by calculation. Compare: about 0.7\,\% for fusion, 100\,\% for annihilation.

\begin{toolkit}{Energy units: J, eV, MeV and u}
$\qty{1}{\eV} = \qty{1.602e-19}{J}$ (data sheet), so $\qty{1}{\mega\eV} = \qty{1.602e-13}{J}$.
Joules to eV: divide by \num{1.602e-19}. Atomic masses are often in atomic mass units:
$\qty{1}{\amu} = \qty{1.661e-27}{kg}$, equivalent to \qty{931.5}{\mega\eV} (data sheet).
Mass in u: multiply $\Delta m$ by 931.5 to get MeV. Mass in kg: use $E = mc^2$ to get J. Do not
mix the two routes in one answer.
\end{toolkit}

\begin{formulacard}{Mass--energy equivalence}
  \cardformula{E = mc^2}
  \begin{symbols}
    \sym{$E$}{energy equivalent to the mass (energy released)}{\si{J}}
    \sym{$m$}{mass; use $\Delta m$, the mass lost, for energy released}{\si{kg}}
    \sym{$c$}{speed of light in a vacuum, \qty{3.00e8}{m.s^{-1}}}{\si{m.s^{-1}}}
  \end{symbols}
  \inwords{multiply a mass by $c^2$ to get the energy it is equivalent to.}
  \whentouse{energy released when mass is lost (fusion, fission, annihilation, burning); the rest
    energy of a particle.}
  \whennottouse{for the kinetic energy of an everyday object ($K = \frac{1}{2}mv^2$), or with a
    mass in g or u.}
  \rearranged{$m = \dfrac{E}{c^2}$ \qquad $\Delta E = \Delta m\,c^2$}
  \tensecond{$\Delta m = \qty{2.0e-6}{kg}$:
    $E = 2.0\times10^{-6} \times (3.00\times10^{8})^2 = \qty{1.8e11}{J}$.}
  \trap{square $c$ ($c^2 = \qty{9.00e16}{m^2.s^{-2}}$). Use only the mass that is \emph{lost}
    (in annihilation, all of it).}
\end{formulacard}

\begin{formulacard}[combine={$P = \frac{\Delta E}{\Delta t}$ and $E = mc^2$}]{Mass lost each second}
  \cardformula{\frac{\Delta m}{\Delta t} = \frac{P}{c^2}}
  \begin{symbols}
    \sym{$\dfrac{\Delta m}{\Delta t}$}{mass lost each second}{\si{kg.s^{-1}}}
    \sym{$P$}{power output: energy given out each second}{\si{W} (\si{J.s^{-1}})}
    \sym{$c$}{speed of light in a vacuum, \qty{3.00e8}{m.s^{-1}}}{\si{m.s^{-1}}}
  \end{symbols}
  \inwords{anything that gives out energy at a steady power $P$ loses mass at the rate
    $\dfrac{P}{c^2}$.}
  \whentouse{the Sun or another star; a power station or reactor with a known power output.}
  \whennottouse{for one single reaction (use $\Delta E = \Delta m\,c^2$ instead).}
  \rearranged{$P = \dfrac{\Delta m}{\Delta t}\,c^2$}
  \tensecond{$P = \qty{1.8e9}{W}$:
    $\dfrac{\Delta m}{\Delta t} = \dfrac{1.8\times10^{9}}{9.00\times10^{16}} = \qty{2.0e-8}{kg.s^{-1}}$.}
  \derive{$P = \dfrac{\Delta E}{\Delta t}$ and $\Delta E = \Delta m\,c^2$, so
    $P = \dfrac{\Delta m\,c^2}{\Delta t}$. Divide both sides by $c^2$.}
  \trap{change MW to W first ($\qty{1}{MW} = \qty{e6}{W}$), and square $c$.}
\end{formulacard}

\begin{workedexample}{Annihilation in a PET scan}\label{ch4:we:annihilation}
In a PET scan, a positron meets an electron. Both are almost at rest, and they annihilate into two
gamma-ray photons. Find (a) the total energy released (b) the energy of each photon in J and MeV
(c) the wavelength of each photon. (d) A scanner detects \num{2.0e8} annihilations each second.
What power do the gamma rays carry?
\given{each mass $m_\mathrm{e} = \qty{9.109e-31}{kg}$ (a positron has the electron's mass);
  $h = \qty{6.626e-34}{J.s}$; $\qty{1}{\mega\eV} = \qty{1.602e-13}{J}$; \num{2.0e8} pairs
  each second}
\find{$\Delta E$, the energy of each photon, its wavelength $\lambda$, and the power $P$}
\formula{$E = mc^2$ with $\Delta m = 2m_\mathrm{e}$; then $E = hf$ and $c = f\lambda$ for one
  photon; $P = \dfrac{\Delta E}{\Delta t}$}
\substitute{\eqlines{\text{(a) } \Delta m &= 2 \times \qty{9.109e-31}{kg} = \qty{1.8218e-30}{kg}\\
  \Delta E &= \qty{1.8218e-30}{kg} \times (\qty{3.00e8}{m.s^{-1}})^2\\
  &= \qty{1.640e-13}{J}\\
  \text{(b) } E_\mathrm{photon} &= \frac{\qty{1.640e-13}{J}}{2} = \qty{8.20e-14}{J}\\
  &= \frac{\num{8.20e-14}}{\num{1.602e-13}}\,\si{\mega\eV} = \qty{0.512}{\mega\eV}\\
  \text{(c) } f &= \frac{E}{h} = \frac{\qty{8.20e-14}{J}}{\qty{6.626e-34}{J.s}} = \qty{1.24e20}{Hz}\\
  \lambda &= \frac{c}{f} = \frac{\qty{3.00e8}{m.s^{-1}}}{\qty{1.24e20}{Hz}} = \qty{2.42e-12}{m}\\
  \text{(d) } P &= \num{2.0e8} \times \qty{1.640e-13}{J} \text{ each second}\\
  &= \qty{3.3e-5}{W}}}
\answer{(a) \final{\qty{1.64e-13}{J}} \quad (b) \final{\qty{8.20e-14}{J} = \qty{0.512}{\mega\eV}}
  each \quad (c)~\final{\qty{2.42e-12}{m}} \quad (d)~\final{\qty{3.3e-5}{W}}}
\checkstep{a wavelength near \qty{e-12}{m} is in the gamma-ray range \cmark. The power is tiny,
  because each photon carries very little energy \cmark. (The accurate photon energy is
  \qty{0.511}{\mega\eV}; the data-sheet $c$ is rounded.)}
\end{workedexample}

\begin{workedexample}{How the Sun loses mass}\label{ch4:we:sun}
The Sun gives out \qty{3.85e26}{W}. (a) How much mass does it lose each second? (b) In the
proton--proton chain, four hydrogen-1 atoms (\qty{1.007825}{\amu} each) end up as one helium-4
atom (\qty{4.002603}{\amu}). Find the energy released each time. (c) How many helium-4 nuclei does
the Sun make each second?
\given{$P = \qty{3.85e26}{W}$ (joules each second); masses in u; $\qty{1}{\amu} =
  \qty{931.5}{\mega\eV}/c^2$; $\qty{1}{\mega\eV} = \qty{1.602e-13}{J}$}
\find{(a) $\dfrac{\Delta m}{\Delta t}$ \quad (b) energy per helium-4 \quad (c) number per second}
\formula{$P = \dfrac{\Delta E}{\Delta t}$ and $\Delta E = \Delta m\,c^2$ give
  $\dfrac{\Delta m}{\Delta t} = \dfrac{P}{c^2}$; \ $E\,(\mathrm{MeV}) = \Delta m\,(\mathrm{u}) \times 931.5$}
\substitute{\eqlines{\text{(a) } \frac{\Delta m}{\Delta t} &= \frac{\qty{3.85e26}{J.s^{-1}}}
    {(\qty{3.00e8}{m.s^{-1}})^2} = \qty{4.28e9}{kg.s^{-1}}\\
  \text{(b) } \Delta m &= 4 \times \qty{1.007825}{\amu} - \qty{4.002603}{\amu} = \qty{0.028697}{\amu}\\
  E &= 0.028697 \times \qty{931.5}{\mega\eV} = \qty{26.73}{\mega\eV}\\
  E &= 26.73 \times \qty{1.602e-13}{J} = \qty{4.282e-12}{J}\\
  \text{(c) } N &= \frac{\qty{3.85e26}{J.s^{-1}}}{\qty{4.282e-12}{J}} = \qty{8.99e37}{s^{-1}}}}
\answer{(a) \final{\qty{4.28e9}{kg}} each second \quad (b) \final{\qty{26.7}{\mega\eV}
  = \qty{4.28e-12}{J}} \quad (c) \final{\num{8.99e37}} each second}
\checkstep{the starting mass is $4 \times \qty{1.007825}{\amu} = \qty{4.0313}{\amu}$, and
  $0.028697 \div 4.0313 = 0.0071$: 0.7\,\% of the mass, as expected for fusion \cmark.
  Even at 4 million tonnes a second, the Sun loses under 0.1\,\% of its mass in 10 billion years.}
\end{workedexample}

%% Added by the Module 7 coverage audit (M7.IQ4.04, burning conventional fuel).
\begin{workedexample}{Burning natural gas}\label{ch4:we:fuel}
A gas heater burns \qty{1.00}{kg} of methane (natural gas) with \qty{3.99}{kg} of oxygen. The
reaction releases \qty{5.55e7}{J} of heat. (a) Find the mass that becomes energy. (b) What
fraction of the starting mass (methane plus oxygen) is this? (c) A good school balance shows
changes as small as \qty{0.01}{g}. Explain why this mass change cannot be measured by weighing.
\given{$\Delta E = \qty{5.55e7}{J}$; \ starting mass $m = 1.00 + 3.99 = \qty{4.99}{kg}$; \
  $c = \qty{3.00e8}{m.s^{-1}}$; \ smallest change the balance shows: $\qty{0.01}{g} =
  \qty{1e-5}{kg}$}
\find{(a) the mass lost $\Delta m$ \quad (b) the fraction $\dfrac{\Delta m}{m}$ \quad (c) how
  $\Delta m$ compares with \qty{1e-5}{kg}}
\formula{$\Delta E = \Delta m\,c^2$. Divide both sides by $c^2$: $\Delta m = \dfrac{\Delta E}{c^2}$}
\substitute{square $c$ first:\\
  \eqlines{c^2 &= (\qty{3.00e8}{m.s^{-1}})^2 = \qty{9.00e16}{m^2.s^{-2}}\\
  \text{(a) } \Delta m &= \frac{\qty{5.55e7}{J}}{\qty{9.00e16}{m^2.s^{-2}}} = \qty{6.17e-10}{kg}\\
  \text{(b) } \frac{\Delta m}{m} &= \frac{\qty{6.17e-10}{kg}}{\qty{4.99}{kg}} = \num{1.24e-10}}\\[5pt]
  (c) how many times bigger the balance's smallest step is than $\Delta m$:\\
  \eqlines{\frac{\qty{1e-5}{kg}}{\qty{6.17e-10}{kg}} &= \num{1.6e4}}}
\answer{(a) \final{$\Delta m = \qty{6.17e-10}{kg}$} \quad (b) \final{\num{1.24e-10}}, about 1 part
  in 8 billion \quad (c) the smallest change the balance can show is about \num{16000} times
  bigger than $\Delta m$. So the change is far too small to weigh. We know it only by
  calculation with $E = mc^2$.}
\checkstep{unit: $\unit{J} \div \unit{m^2.s^{-2}} = \unit{kg.m^2.s^{-2}} \div \unit{m^2.s^{-2}} =
  \unit{kg}$ \cmark. The size is close to the petrol example above (\qty{5.1e-10}{kg} for
  \qty{1.00}{kg} of fuel) \cmark. Fusion turns 0.7\,\% of its mass into energy: more than 50
  million times the fraction here.}
\end{workedexample}

\begin{mistakes}
\begin{itemize}
  \item \xmark\ One photon from annihilation. \cmark\ Two photons in opposite directions
        (momentum is conserved), and $\Delta m = 2m_\mathrm{e}$.
  \item \xmark\ Multiplying a mass in u by $c^2$. \cmark\ Mass in u: multiply by 931.5 to get MeV.
  \item \xmark\ ``Only nuclear reactions turn mass into energy.'' \cmark\ In every process that
        releases energy, the products have less mass. In chemistry the difference is too small to
        measure.
\end{itemize}
\end{mistakes}

\begin{examtip}
``Calculate the energy released'' is worth 2--3 marks. Full marks need:
\begin{enumerate}
  \item $\Delta m$ from the correct masses (mass before minus mass after);
  \item $\Delta E = \Delta m\,c^2$ (or $\times\,931.5$ MeV for masses in u);
  \item the answer with its unit, after any unit change.
\end{enumerate}
Annihilation: say \emph{two} photons, and give each one half of the energy. Multiple-choice
questions ask about fuel burning in a sealed jar. Its mass falls very slightly, but only as the
heat and light leave the jar. Questions can also ask about the reverse process: in an
accelerator, the energy of a collision can make new particles that have mass.
\end{examtip}

\begin{quickcheck}
  \item A nuclear reactor releases energy at \qty{3.0e9}{W}. How much mass becomes energy
        (a) each second (b) each year ($\qty{1}{\yr} = \qty{3.16e7}{s}$)?
  \item A proton and an antiproton at rest annihilate. Assume they make two photons (in real
        life they usually make other particles called pions). Find each photon's energy in J
        and in MeV.
  \item Burning \qty{1.0}{kg} of coal releases \qty{3.0e7}{J}. Find the decrease in mass. Why can
        it not be measured?
  \item Deuterium (hydrogen-2, \qty{2.014102}{\amu}) and tritium (hydrogen-3,
        \qty{3.016049}{\amu}) fuse into helium-4
        (\qty{4.002603}{\amu}) and a neutron (\qty{1.008665}{\amu}). Find the energy released in MeV
        and in J.
\end{quickcheck}

%% ==================================================================
\section{Key experiments}\label{ch4:sec:experiments}
\covers{M7.IQ4.01, M7.IQ4.02, M7.IQ4.03}

\begin{refresher}{Validity, reliability and accuracy}
These are the same words as in \cref{ch1:sec:speed}.
\textbf{Validity}: the method really tests what it claims to, with other variables controlled
(a fair test). \textbf{Reliability}: repeats give consistent results. \textbf{Accuracy}: how
close a result is to the true (or predicted) value.
\end{refresher}

\begin{experiment}{Michelson and Morley (1887): looking for the aether}
\aim{To detect the Earth's motion through the aether. An \term{interferometer} is an
  instrument that splits light into two beams and joins them again. Here it compares the speed
  of light in two directions at right angles.}
\apparatus{
  \begin{sidebyside}[0.52]
    \centering
    \begin{tikzpicture}[x=0.6cm, y=0.6cm]
      \fill[hscmuted!12, rounded corners=2mm] (-0.9,-2.9) rectangle (6.9,3.5);
      \node[diagramlabel, font=\scriptsize, text=hscmuted, anchor=north west, align=left] at (-0.85,3.45)
        {stone slab\\[-2pt]on mercury};
      \draw[object] (-0.7,-0.3) rectangle (0,0.3);
      \node[diagramlabel, font=\scriptsize, above] at (-0.35,0.3) {source};
      \draw[ray] (0,0.06) -- (2.95,0.06);
      \draw[hscink, line width=1.6pt] (2.65,-0.35) -- (3.35,0.35);
      \node[diagramlabel, font=\scriptsize, anchor=north east, align=right] at (2.8,-0.35)
        {half-silvered\\[-2pt] mirror};
      \draw[ray] (3.05,0.06) -- (5.9,0.06);
      \draw[ray] (5.9,-0.08) -- (3.08,-0.08);
      \draw[hscink, line width=2pt] (6.0,-0.5) -- (6.0,0.5);
      \node[diagramlabel, font=\scriptsize, anchor=south] at (6.0,0.55) {mirror};
      \draw[ray] (2.94,0.1) -- (2.94,2.9);
      \draw[ray] (3.08,2.9) -- (3.08,0.1);
      \draw[hscink, line width=2pt] (2.5,3.0) -- (3.5,3.0);
      \node[diagramlabel, font=\scriptsize, anchor=west] at (3.55,3.0) {mirror};
      \draw[ray] (3.0,-0.12) -- (3.0,-1.75);
      \draw[object] (2.65,-2.35) rectangle (3.35,-1.8);
      \node[diagramlabel, font=\scriptsize, anchor=west] at (3.4,-2.05) {telescope};
      \node[diagramlabel, font=\scriptsize, text=hscmuted] at (4.5,0.45) {arm 1};
      \node[diagramlabel, font=\scriptsize, text=hscmuted, anchor=east] at (2.85,1.7) {arm 2};
    \end{tikzpicture}
    \captionof{figure}{The interferometer, seen from above.}\label{ch4:fig:mm}
  \nextside
    A light source; a half-silvered mirror that splits the beam; two equal arms at right angles
    with mirrors (each light path about \qty{11}{m}); a telescope to watch the interference
    fringes; a heavy stone slab floating on mercury, so the apparatus could turn smoothly.
  \end{sidebyside}}
\method{\begin{enumerate}
  \item Split one beam into two at right angles. Each goes along its arm and back, and the two
        recombine to make interference fringes (bright and dark bands).
  \item Slowly rotate the apparatus and watch for a fringe shift. Repeat at different times of
        day and year, as the Earth's motion changes direction.
\end{enumerate}}
\results{The aether idea predicted a shift of about 0.4 of a fringe. The shift seen was less than
  0.02 of a fringe: a null result (the expected effect was not found).}
\conclusion{No aether wind: light has the same speed in all directions on the moving Earth.}
\significance{It found none of the effects the aether predicted, so there was no evidence for an
  aether. The result agrees with both postulates.}
\validity{\begin{itemize}
  \item \textbf{Validity:} both beams came from one source and were compared with each other,
        so only turning the apparatus (changing direction) could shift the fringes.
  \item \textbf{Reliability:} repeated many times; modern laser versions always give a null result.
  \item \textbf{Accuracy:} it could detect 0.01 fringe, far less than the 0.4 expected.
\end{itemize}}
\end{experiment}

\begin{experiment}{Cosmic-ray muons on a mountain (Frisch and Smith, 1963)}
\aim{To test time dilation by counting fast muons on a mountain top and at sea level.}
\apparatus{
  \begin{sidebyside}[0.4]
    \centering
    \begin{tikzpicture}[x=0.6cm, y=0.6cm]
      \fill[hscmuted!25] (0,0) -- (2.4,2.6) -- (4.6,0) -- cycle;
      \draw[hscink, line width=0.8pt] (-0.4,0) -- (6.6,0);
      \draw[object] (2.1,2.6) rectangle (2.7,3.0);
      \draw[object] (5.3,0) rectangle (5.9,0.4);
      \foreach \x in {2.25,2.55} \draw[velocity, line width=0.9pt] (\x,4.2) -- (\x,3.1);
      \foreach \x in {5.45,5.75} \draw[velocity, line width=0.9pt] (\x,4.2) -- (\x,0.5);
      \node[diagramlabel, font=\scriptsize, anchor=south] at (4.0,4.2) {muons};
      \draw[<->, hscmuted, line width=0.6pt] (1.0,0) -- (1.0,2.6);
      \node[diagramlabel, font=\scriptsize, anchor=east] at (0.95,1.3) {\qty{1.9}{km}};
      \node[diagramlabel, font=\scriptsize, anchor=north] at (5.6,-0.05) {sea level};
    \end{tikzpicture}
    \captionof{figure}{Counting muons at two heights.}\label{ch4:fig:mountain}
  \nextside
    A scintillator (it flashes when a muon passes) under a block of iron. Only muons with
    speeds near $0.995c$ stopped inside it and were counted. The same detector was used on
    Mount Washington (USA), about \qty{1.9}{km} up, and near sea level. At sea level some iron
    was removed, to make up for the energy muons lose in \qty{1.9}{km} of air. So the same
    group of muons was counted.
  \end{sidebyside}}
\method{\begin{enumerate}
  \item Count the muons near $0.995c$ arriving each hour on the mountain, then at sea level.
  \item Time each stopped muon from the flash when it arrives to the flash when it decays. The
        average gives the mean lifetime at rest (\qty{2.2}{\micro\second}).
\end{enumerate}}
\results{About 563 muons per hour arrived on top, and 408 per hour at sea level. The trip down
  takes about \qty{6.4}{\micro\second} in the Earth frame. That is nearly 3 lifetimes. Without
  time dilation, only about 30 per hour should survive.}
\conclusion{Far more muons survive than classical physics predicts. The counts match a
  time-dilation factor of about 9. Relativity predicts this value once you allow for the muons
  slowing down in the air.}
\significance{Direct evidence for time dilation (and, in the muon frame, length contraction).}
\validity{\begin{itemize}
  \item \textbf{Validity:} one narrow speed range, the same detector, and the iron adjusted for
        the air, so the same muons were compared.
  \item \textbf{Reliability:} hundreds of counts per hour for many hours; repeated by others.
  \item \textbf{Accuracy:} decay is random and muons slow in air, so averages and corrections are
        needed. The result agreed with the prediction within its uncertainty.
\end{itemize}}
\end{experiment}

\begin{experiment}{Atomic clocks on aeroplanes (Hafele and Keating, 1971)}
\aim{To test time dilation with very accurate atomic clocks flown around the world.}
\apparatus{
  \begin{sidebyside}[0.45]
    \centering
    \begin{tikzpicture}[x=0.8cm, y=0.8cm]
      \fill[hscvel!10] (0,0) circle (1.3);
      \draw[hscvel!60, line width=0.8pt] (0,0) circle (1.3);
      \fill[hscink] (0,0) circle (0.05);
      \node[diagramlabel, font=\scriptsize, below] at (0,-0.02) {North Pole};
      \draw[->, hscmuted, line width=0.8pt] (40:0.6) arc (40:140:0.6);
      \node[diagramlabel, font=\scriptsize, text=hscmuted] at (0,0.85) {Earth spins};
      \fill[hscink] (270:1.3) circle (0.07);
      \draw[velocity] (270:1.3) -- ++(0:0.9);
      \node[diagramlabel, font=\scriptsize, anchor=north] at (270:1.4) {ground clock};
      \fill[hscforce] (330:1.55) circle (0.07);
      \draw[velocity] (330:1.55) -- ++(60:1.3);
      \node[diagramlabel, font=\scriptsize, anchor=south, align=center] at (1.75,1.55)
        {eastward plane:\\[-2pt] fastest};
      \fill[hscacc] (150:1.55) circle (0.07);
      \draw[velocity] (150:1.55) -- ++(240:0.42);
      \node[diagramlabel, font=\scriptsize, anchor=south, align=center] at (-1.6,1.0)
        {westward plane:\\[-2pt] slowest};
    \end{tikzpicture}
    \captionof{figure}{Speeds relative to the Earth's centre, seen from above the North Pole
      (not to scale). Arrow length shows speed.}\label{ch4:fig:hk}
  \nextside
    Four caesium \term{atomic clocks} (very accurate clocks that use the vibrations of atoms);
    normal passenger jets; reference atomic clocks on the ground (US Naval Observatory); flight
    records of speed, height and route.
  \end{sidebyside}}
\method{Use the Earth's centre as an (almost) inertial frame (\cref{ch4:fig:hk}). The ground
  moves east as the Earth spins. The eastward plane moves fastest, so its clock should run
  slowest. The westward plane partly cancels the spin, so it moves slowest and its clock should
  run fastest.
  \begin{enumerate}
  \item Synchronise the four clocks with the ground clocks.
  \item Fly them around the world eastward, then westward. Compare with the ground clocks after
        each trip.
\end{enumerate}}
\results{Time gained (+) or lost (\textminus) by the flying clocks, in nanoseconds:
  \begin{center}\small
  \begin{tabular}{@{}L{18mm}C{25mm}C{25mm}C{20mm}C{18mm}@{}}
    \toprule
    \colhead{Trip} & \colhead{Predicted: speed (special relativity)} & \colhead{Predicted: height
      (gravity)} & \colhead{Predicted total} & \colhead{Measured}\\
    \midrule
    Eastward & $-184 \pm 18$ & $+144 \pm 14$ & $-40 \pm 23$ & $-59 \pm 10$\\
    Westward & $+96 \pm 10$ & $+179 \pm 18$ & $+275 \pm 21$ & $+273 \pm 7$\\
    \bottomrule
  \end{tabular}
  \end{center}
  The ``height'' effect (higher clocks run faster) comes from general relativity, Einstein's
  later theory of gravity. It is not in the HSC course. The special-relativity effects alone
  ($-184$ and $+96$) do not match the measurements; adding the height effect does.}
\conclusion{Measured and predicted times agree within the uncertainties: moving clocks run slow,
  even at aeroplane speeds.}
\significance{The first test with man-made atomic clocks carried on aircraft. GPS (satellite
  navigation) satellites use the same corrections (about
  \qty{38}{\micro\second} a day); without them, positions would drift about \qty{11}{km} a day.}
\validity{\begin{itemize}
  \item \textbf{Validity:} the same ground clocks were the reference; the gravity effect was included.
  \item \textbf{Reliability:} only one trip each way, on uncontrolled flights; later repeats agreed.
  \item \textbf{Accuracy:} the effects are only tens to hundreds of nanoseconds, so four clocks were
        averaged; both results agree within their uncertainties.
\end{itemize}}
\end{experiment}

\begin{toolkit}{When $v$ is tiny compared with $c$}
For an aeroplane, $\dfrac{v^2}{c^2}$ is about $10^{-12}$. A calculator then shows
$\sqrt{1 - v^2/c^2}$ as exactly 1, so the time difference seems to be zero. Use this shortcut
instead:
\[ t - t_0 \approx t \times \frac{v^2}{2c^2} \]
{\centering\hscbadgenotonsheet\par}
Where it comes from:
\begin{itemize}
  \item For a tiny number $x$, $\sqrt{1 - x} \approx 1 - \dfrac{x}{2}$. Check:
        $\left(1 - \dfrac{x}{2}\right)^2 = 1 - x + \dfrac{x^2}{4}$, and $\dfrac{x^2}{4}$ is far
        too small to matter.
  \item So $t_0 = t\sqrt{1 - x} \approx t - t \times \dfrac{x}{2}$.
  \item Rearrange: $t - t_0 \approx t \times \dfrac{x}{2}$. Then put in $x = \dfrac{v^2}{c^2}$.
\end{itemize}
Example: a plane flies at \qty{200}{m.s^{-1}} for \qty{10.0}{h} ($t = \qty{36000}{s}$).
Then $\dfrac{v^2}{c^2} = \left(\dfrac{200}{\num{3.00e8}}\right)^2 = \num{4.44e-13}$, and half
of this is $\dfrac{v^2}{2c^2} = \num{2.22e-13}$. So
$t - t_0 = \num{36000} \times \num{2.22e-13} = \qty{8.0e-9}{s} = \qty{8.0}{ns}$.
\end{toolkit}

%% Added by the Module 7 coverage audit (M7.IQ4.02, atomic clocks analysed quantitatively).
\begin{workedexample}{Clocks flown east and west}\label{ch4:we:clocks}
A jet flies along the equator for \qty{40.0}{h} at a steady \qty{240}{m.s^{-1}} relative to
the ground. One trip goes east and another goes west. The Earth turns once in \qty{24.0}{h}. Use the Earth's centre as the inertial frame, and use special relativity only
(ignore the height effect). After each trip, how far ahead of, or behind, a clock on the ground
at the equator is the flying clock? (a) Eastward trip. (b) Westward trip.
\given{$r_\mathrm{E} = \qty{6.371e6}{m}$ (data sheet); \ $T = 24.0 \times \qty{3600}{s} =
  \qty{86400}{s}$; \ jet speed over the ground $v_\mathrm{jet} = \qty{240}{m.s^{-1}}$; \
  $t = 40.0 \times \qty{3600}{s} = \qty{1.44e5}{s}$}
\find{the speed of each clock relative to the Earth's centre; how far each clock falls behind;
  then flying clock minus ground clock}
\formula{$v = \dfrac{2\pi r}{T}$ (on the sheet) for the ground clock. The eastward jet moves
  with the spin: $v_\mathrm{ground} + v_\mathrm{jet}$. The westward jet moves against it:
  $v_\mathrm{ground} - v_\mathrm{jet}$. Then the toolkit shortcut $t - t_0 \approx t \times
  \dfrac{v^2}{2c^2}$ for each clock.}
\substitute{speeds relative to the Earth's centre:\\
  \eqlines{v_\mathrm{ground} &= \frac{2\pi \times \qty{6.371e6}{m}}{\qty{86400}{s}}
    = \qty{463.3}{m.s^{-1}}\\
  v_\mathrm{east} &= \qty{463.3}{m.s^{-1}} + \qty{240}{m.s^{-1}} = \qty{703.3}{m.s^{-1}}\\
  v_\mathrm{west} &= \qty{463.3}{m.s^{-1}} - \qty{240}{m.s^{-1}} = \qty{223.3}{m.s^{-1}}}\\[5pt]
  how far each clock falls behind a clock at rest at the Earth's centre
  ($2c^2 = \qty{1.80e17}{m^2.s^{-2}}$; \ $\qty{1}{ns} = \qty{e-9}{s}$). Ground clock:\\
  \eqlines{\frac{v^2}{2c^2} &= \frac{(\qty{463.3}{m.s^{-1}})^2}{\qty{1.80e17}{m^2.s^{-2}}}
    = \num{1.1925e-12}\\
  t - t_0 &= \qty{1.44e5}{s} \times \num{1.1925e-12} = \qty{1.717e-7}{s} = \qty{171.7}{ns}}\\
  Eastward clock:\\
  \eqlines{\frac{v^2}{2c^2} &= \frac{(\qty{703.3}{m.s^{-1}})^2}{\qty{1.80e17}{m^2.s^{-2}}}
    = \num{2.7479e-12}\\
  t - t_0 &= \qty{1.44e5}{s} \times \num{2.7479e-12} = \qty{3.957e-7}{s} = \qty{395.7}{ns}}\\
  Westward clock:\\
  \eqlines{\frac{v^2}{2c^2} &= \frac{(\qty{223.3}{m.s^{-1}})^2}{\qty{1.80e17}{m^2.s^{-2}}}
    = \num{2.7702e-13}\\
  t - t_0 &= \qty{1.44e5}{s} \times \num{2.7702e-13} = \qty{3.99e-8}{s} = \qty{39.9}{ns}}\\[5pt]
  compare with the ground clock. Each clock reads $t$ minus its own lag, so
  (flying clock) $-$ (ground clock) $=$ (ground lag) $-$ (flying lag):\\
  \eqlines{\text{(a) } \qty{171.7}{ns} - \qty{395.7}{ns} &= \qty{-224.0}{ns}\\
  \text{(b) } \qty{171.7}{ns} - \qty{39.9}{ns} &= +\qty{131.8}{ns}}}
\answer{(a) the eastward clock is \final{\qty{224}{ns} behind} the ground clock. (b) The
  westward clock is \final{\qty{132}{ns} ahead} of it. (3 significant figures.)}
\checkstep{the signs match the speed predictions in the Hafele--Keating table ($-184$ and
  $+96$\,ns) \cmark. The real flights were not along the equator and took different times, so
  the numbers differ. The eastward clock moves fastest, so it runs slowest \cmark.}
\end{workedexample}

\begin{experiment}{Evidence from particle accelerators}
\aim{To test the speed limit, relativistic momentum and time dilation with very fast particles.}
\apparatus{
  \begin{figure}
    \centering
    \begin{tikzpicture}[x=1cm, y=1cm]
      \draw[object, fill=hscbrand!5] (0,0) rectangle (2.3,0.8);
      \node[diagramlabel, font=\scriptsize, align=center] at (1.15,0.4)
        {accelerator\\[-2pt](electric field)};
      \draw[velocity] (2.35,0.4) -- (3.2,0.4);
      \node[diagramlabel, font=\scriptsize, text=hscvel, anchor=south] at (5.9,0.45) {electrons};
      \fill[hscmuted] (3.3,0.05) rectangle (3.48,0.75);
      \node[diagramlabel, font=\scriptsize, anchor=north, align=center] at (3.39,0)
        {detector A\\[-2pt](start timer)};
      \draw[hscvel, dashed, line width=0.8pt, ->] (3.5,0.4) -- (8.25,0.4);
      \fill[hscmuted!60] (8.3,-0.05) rectangle (8.9,0.85);
      \draw[hscink, line width=0.8pt] (8.6,0.85) -- (8.6,1.35);
      \fill[hscforce] (8.6,0.85) circle (0.07);
      \node[diagramlabel, font=\scriptsize, anchor=west] at (8.7,1.25) {thermometer};
      \node[diagramlabel, font=\scriptsize, anchor=north, align=center] at (8.6,-0.1)
        {aluminium target\\[-2pt](stop timer)};
      \draw[<->, hscmuted, line width=0.6pt] (3.39,1.0) -- node[above, text=hscink,
        font=\scriptsize] {flight path \qty{8.4}{m}} (8.3,1.0);
    \end{tikzpicture}
    \caption{Bertozzi's apparatus (1964), simplified. The electrons are timed between
      detector A and the target; the target's temperature rise gives their
      energy.}\label{ch4:fig:bertozzi-app}
  \end{figure}
  Bertozzi (1964): an accelerator (electric fields speed up the electrons); a detector and an
  aluminium target \qty{8.4}{m} apart to time them; a thermometer on the target. CERN (1977): a
  ring of magnets that stores fast muons, and detectors that record when each muon decays.}
\method{\begin{enumerate}
  \item Bertozzi (1964): give electrons \qtyrange{0.5}{15}{\mega\eV} of kinetic energy; time them
        over the flight path; measure their energy separately (heat in the target).
  \item CERN (1977): store muons in the magnetic ring, moving at $0.999418c$. At this speed the
        time-dilation factor $\dfrac{1}{\sqrt{1 - v^2/c^2}}$ is 29.3. (Keep all six digits. Near
        $c$ the factor changes fast: $0.9994c$ would give 28.9.) Measure how long the muons live.
\end{enumerate}}
\results{
  \begin{sidebyside}[0.5]
    \centering
    \begin{tikzpicture}
      \begin{axis}[hsc, width=0.95\linewidth, height=0.5\linewidth,
          xmin=0, xmax=15, ymin=0, ymax=2, ytick={0,1,2},
          xlabel={kinetic energy (\si{\mega\eV})}, ylabel={speed $v/c$}]
        \addplot[hscseriesA, no marks, line width=1.2pt, domain=0:15, samples=150]
          {sqrt(1-1/(1+x/0.511)^2)};
        \addplot[hscseriesB, no marks, dashed, line width=1.2pt, domain=0:1.02, samples=60]
          {sqrt(2*x/0.511)};
        \draw[guide] (axis cs:0,1) -- (axis cs:15,1);
        \node[font=\scriptsize, anchor=south east] at (axis cs:15,1.03) {speed of light};
        \node[font=\scriptsize, anchor=north east, text=hscseriesA, align=right] at (axis cs:15,0.93)
          {relativity (matches\\ the measurements)};
        \node[font=\scriptsize, anchor=west, text=hscseriesB] at (axis cs:1.2,1.75)
          {classical $K = \frac{1}{2}mv^2$};
      \end{axis}
    \end{tikzpicture}
    \captionof{figure}{Electron speed against kinetic energy. The dashed line is the classical
      prediction; real electrons follow the solid line.}\label{ch4:fig:bertozzi}
  \nextside
    The electrons' speed approached $c$ but never passed it (\cref{ch4:fig:bertozzi}). At
    \qty{15}{\mega\eV}, $\frac{1}{2}mv^2$ predicts $7.7c$; the measured speed was just below $c$.

    The stored muons lived \qty{64.4}{\micro\second}. That is 29.3 times their lifetime at rest
    (\qty{2.197}{\micro\second}), as time dilation predicts, to about 0.1\,\%. Check:
    $\qty{2.197}{\micro\second} \times 29.3 = \qty{64.4}{\micro\second}$.
  \end{sidebyside}}
\conclusion{The speed limit $c$ and time dilation are confirmed. Steering magnets must match $p_\mathrm{v}$.}
\significance{Every accelerator (such as the LHC) is designed with $p_\mathrm{v}$: a daily test.}
\validity{\begin{itemize}
  \item \textbf{Validity:} speed and energy were measured separately: a fair comparison.
  \item \textbf{Reliability:} huge numbers of particles and many repeated runs.
  \item \textbf{Accuracy:} precise speeds and lifetimes that agree with relativity.
\end{itemize}}
\end{experiment}

\begin{experiment}{Evidence from cosmology: slow-motion supernovae}
\aim{To test whether events in very distant galaxies are measured to run slow.}
\apparatus{
  \begin{figure}
    \centering
    \begin{tikzpicture}[x=1cm, y=1cm]
      \fill[hscmuted!25] (0,0) ellipse (0.85 and 0.4);
      \fill[hscacc] (0.25,0.08) circle (0.11);
      \node[diagramlabel, font=\scriptsize, anchor=north, align=center] at (0,-0.45)
        {distant galaxy\\[-2pt]with a supernova};
      \draw[hscray, line width=1.1pt, domain=0:3.25, samples=240, smooth]
        plot ({1.0+\x}, {0.12*sin(deg(34.91*ln(1+0.9*\x)))});
      \draw[hscray, line width=1.1pt, ->] (4.25,0) -- (4.45,0);
      \node[diagramlabel, font=\scriptsize, anchor=north, align=center, text=hscray] at (2.7,-0.25)
        {light, stretched by the\\[-2pt]expanding Universe};
      \draw[object, fill=hscbrand!5] (4.5,-0.3) -- (6.0,-0.18) -- (6.0,0.18) -- (4.5,0.3) -- cycle;
      \node[diagramlabel, font=\scriptsize, anchor=south] at (5.25,0.32) {telescope};
      \draw[hscmuted, line width=0.8pt, ->] (6.05,0.05) -- (6.9,0.55);
      \draw[hscmuted, line width=0.8pt, ->] (6.05,-0.05) -- (6.9,-0.55);
      \draw[object, fill=hscbrand!5] (6.95,0.35) rectangle (7.55,0.8);
      \node[diagramlabel, font=\scriptsize, anchor=west, align=left] at (7.6,0.58)
        {camera: brightness\\[-2pt]night after night};
      \draw[object, fill=hscbrand!5] (6.95,-0.8) rectangle (7.55,-0.35);
      \node[diagramlabel, font=\scriptsize, anchor=west, align=left] at (7.6,-0.58)
        {spectroscope: spectrum\\[-2pt]$\to$ redshift $z$};
    \end{tikzpicture}
    \caption{Measuring a distant supernova, simplified.}\label{ch4:fig:sn-app}
  \end{figure}
  Telescopes with cameras that measure a supernova's brightness night after night;
  \term{spectroscopes} (instruments that spread light into a spectrum, \cref{ch1:sec:spectra})
  that measure its
  \term{redshift} $z$ (how much its light is stretched by the expanding Universe).}
\method{A \term{Type Ia supernova} is an exploding star. They all brighten and fade in nearly the
  same way over a few weeks. Measure this ``light curve'' and $z$ for hundreds of supernovae at
  different distances.}
\results{
  \begin{sidebyside}[0.5]
    \centering
    \begin{tikzpicture}
      \begin{axis}[hsc, width=0.95\linewidth, height=0.56\linewidth,
          xmin=0, xmax=120, ymin=0, ymax=1.2, ytick={0,0.5,1},
          yticklabel style={/pgf/number format/.cd, fixed, precision=1},
          xlabel={time seen from Earth (days)}, ylabel={relative brightness}]
        \addplot[hscseriesA, no marks, line width=1.2pt, domain=0:120, samples=160]
          {(x/18)^2*exp(2*(1-x/18))};
        \addplot[hscseriesB, no marks, dashed, line width=1.2pt, domain=0:120, samples=160]
          {(x/36)^2*exp(2*(1-x/36))};
        \node[font=\scriptsize, anchor=south west, text=hscseriesA] at (axis cs:21,1.0) {nearby};
        \node[font=\scriptsize, anchor=north east, text=hscseriesB, align=left] at (axis cs:119,1.15)
          {distant ($z = 1$):\\ 2 times slower};
      \end{axis}
    \end{tikzpicture}
    \captionof{figure}{Light curves of two Type Ia supernovae (simplified; each scaled to the
      same peak).}\label{ch4:fig:supernova}
  \nextside
    Distant supernovae brighten and fade more slowly (\cref{ch4:fig:supernova}). One at
    $z = 1$ takes twice as long as a nearby one.
  \end{sidebyside}
  The stretch factor is $1 + z$ \hscbadgenotonsheet. The same factor stretches the wavelength of
  their light:
  \[ 1 + z = \frac{\lambda_\mathrm{observed}}{\lambda_\mathrm{emitted}} \]
  Example:
  a supernova that fades in 20 days nearby takes $20 \times (1 + 0.80) = 36$ days at $z = 0.80$.

  In 2023 an Australian-led team found the same slow motion in quasars (very bright galaxy
  centres). These were seen when the Universe was just over a billion years old.}
\conclusion{Events in distant galaxies are measured to run slow: \term{cosmological time
  dilation}. It comes mainly from the expansion of the Universe (general relativity, which grew
  out of special relativity).}
\significance{Time is not absolute, even across the Universe. It rules out the ``tired light''
  idea (that light just loses energy on its long trip), which predicts no stretch in time.}
\validity{\begin{itemize}[beginpenalty=10000]% layout: no break after the heading
  \item \textbf{Validity:} Type Ia supernovae are very alike, so they are fair comparisons; their
        spectra also age more slowly, an independent check.
  \item \textbf{Reliability:} hundreds of supernovae and several independent teams agree.
  \item \textbf{Accuracy:} single curves are uncertain, but the average stretch matches $1 + z$.
\end{itemize}}
\end{experiment}

\begin{examtip}
``Recount an experiment that confirmed time dilation'' (about 3 marks) needs:
\begin{enumerate}
  \item the experiment named, and what was measured (muon counts at two heights, or atomic
        clocks flown east and west);
  \item the result with numbers (408 muons per hour at sea level instead of about 30; or
        $-59$ and $+273$\,ns);
  \item a link from the result to the prediction of special relativity.
\end{enumerate}
``Assess the validity (or reliability)'' (2--3 marks): name one specific feature of the method
(for example, four clocks were averaged). Say what it controls or repeats, and how far it lets
you trust the result. A definition alone earns no marks.
\end{examtip}

\begin{examtip}
``Assess the role of thought experiments'' is worth about 5 marks. Full marks need:
\begin{enumerate}
  \item what a thought experiment is;
  \item two examples and what each predicted (light clock: time dilation; train and lightning:
        simultaneity);
  \item a strength (it explores speeds near $c$ that could not be tested in 1905) and a
        limitation (it is not evidence);
  \item the experiments that later confirmed them (muons, Hafele--Keating, accelerators), and
        a judgement: essential for predictions, but accepted only after experiments.
\end{enumerate}
``Analyse the consequences of special relativity'' (about 8 marks): link each consequence to
its evidence: time dilation (muons, atomic clocks), length contraction (the muon frame),
momentum and the speed limit (accelerators).
\end{examtip}

\begin{quickcheck}
  \item Outline how counting muons on a mountain and at sea level gives evidence for time
        dilation.
  \item In the Hafele--Keating experiment, explain why the eastward clocks lost time compared
        with the ground clocks.
  \item A jet flies at \qty{250}{m.s^{-1}} (relative to the Earth's centre) for \qty{48.0}{h}.
        Find (a) $v^2/c^2$ (b) how far its clock falls behind. Use the toolkit shortcut.
  \item A distant Type Ia supernova at $z = 0.50$ takes 30 days to fade, measured from Earth. How
        long would the same fading take for a nearby one?
\end{quickcheck}
